% ============================================================
% Chapter 20 --- Change of variables, Fubini and Feynman's trick
% Originally: content/week-08.tex --- \subsection{Change of variables} and
%             \subsection{Fubini's theorem} (both promoted to sections) together with
%             \section{Feynman's trick} (Corsin Week 8), plus problem 8.5 from sheet 8.
% ============================================================
\chapter{Change of Variables, Fubini and Feynman's Trick}
\label{ch:change_of_variables}

% Transition: Claude Opus 5 (Medium)
With the integral defined, the problem becomes computing one, and three tools do nearly all of
that work. Change of variables trades an awkward domain for a convenient one and pays a Jacobian
factor for the trade. Fubini reduces an $n$-dimensional integral to $n$ one-dimensional ones.
Differentiating under the integral sign reaches integrals with no elementary antiderivative at
all, by turning the integral into an ODE in a parameter it never had. A short closing section puts
all three to work at once on the centroid of a region, which is the first quantity in the course
that is defined by an integral rather than computed by one.

% Originally: content/week-08.tex, \subsection{Change of variables} --- promoted to a section

\section{Change of variables}
\label{sec:change_of_variables}

\begin{theorem}[Change of variables]
\label{thm:change_of_variables}
Suppose $U,V \subseteq \mathbb{R}^n$ are open sets and $\Phi : U \to V$ is a
$C^1$-diffeomorphism. If $A \subseteq U$ is Jordan measurable with $\overline{A} \subseteq U$,
and $f : A \to \mathbb{R}$ is continuous, then
\[\int_A f(x)\,dx
  = \int_{\Phi(A)} \frac{f\circ\Phi^{-1}(y)}{\big|\det \Jac\Phi\big(\Phi^{-1}(y)\big)\big|}\,dy.\]
\end{theorem}

\begin{ainote}
Corsin uses \newterm{Jordan measurable} \germanfor{Jordan-messbar} here without defining it;
the definition is in the lecture (ch.~13). For this chapter it is enough to know that boxes,
and the bounded regions cut out by finitely many smooth curves that appear in the examples, are
all Jordan measurable.

On the shape of the formula: it is often quoted the other way round, as
$\int_{\Phi(A)} f = \int_A (f\circ\Phi)\,\lvert\det\Jac\Phi\rvert$. The version above is the same
statement with $\Phi$ pushed to the other side, which is why the determinant appears
\emph{divided} rather than multiplied. Which form is convenient depends on whether it is $\Phi$
or $\Phi^{-1}$ that you can write down. In \cref{ex:change_of_variables_domain} below it is
$\Phi$, so the dividing form is the one that saves work.

% Added 2026-08-13 (Claude Opus 5, Extra): the two halves of the proof are worked out at the very
% end of this section, six hundred lines below, and nothing here said so. A reader has no way to
% find them.
The theorem is not proved here, but the two steps that carry its proof are, and both sit at the
end of this section. \Cref{ex:linear_substitution_elementary_matrices} settles the linear case,
where the determinant enters and everything else is bookkeeping.
\Cref{ex:cube_covering_lemma_substitution} settles the geometric step that a Lipschitz estimate
cannot reach, namely that a diffeomorphism close to the identity cannot let a cube collapse.
\end{ainote}

% Source: Jérôme Paschoud/Notizen/Woche 8.2 Riemann-Integral.pdf, p. 2
% Generator: Gemini 3.1 Pro (High)
\begin{remark}[Intuition for the Jacobian determinant]
From linear algebra, we know that for a linear map given by a matrix $L$, the absolute value of its determinant $\lvert\det L\rvert$ describes how much the map scales volumes. For example, the map $L = \left(\begin{smallmatrix} 2 & 3 \\ 1 & -1 \end{smallmatrix}\right)$ expands areas by a factor of $\lvert\det L\rvert = \lvert-2 - 3\rvert = 5$.

\begin{center}
\begin{tikzpicture}[>=Latex, scale=0.8]
  % Left side: Unit square
  \draw[->, gray] (-0.5, 0) -- (2, 0);
  \draw[->, gray] (0, -0.5) -- (0, 2);
  \draw[thick, blue, fill=blue!10] (0,0) -- (1,0) -- (1,1) -- (0,1) -- cycle;
  \draw[->, thick, blue] (0,0) -- (1,0) node[below right, font=\scriptsize, text=black] {$\left(\begin{smallmatrix} 1 \\ 0 \end{smallmatrix}\right)$};
  \draw[->, thick, blue] (0,0) -- (0,1) node[above left, font=\scriptsize, text=black] {$\left(\begin{smallmatrix} 0 \\ 1 \end{smallmatrix}\right)$};
  
  % Arrow in middle
  \draw[->, thick] (2.5, 0.5) -- (4, 0.5) node[midway, above] {$L$};
  
  % Right side: Parallelogram
  \begin{scope}[shift={(6, 1)}]
    \draw[->, gray] (-0.5, 0) -- (5.5, 0);
    \draw[->, gray] (0, -2.5) -- (0, 2.5);
    
    \coordinate (O) at (0,0);
    \coordinate (V1) at (2,1);
    \coordinate (V2) at (3,-1);
    \coordinate (V3) at (5,0);
    
    \draw[thick, blue, fill=blue!10] (O) -- (V1) -- (V3) -- (V2) -- cycle;
    \draw[->, thick, blue] (O) -- (V1) node[above left, font=\scriptsize, text=black] {$\left(\begin{smallmatrix} 2 \\ 1 \end{smallmatrix}\right)$};
    \draw[->, thick, blue] (O) -- (V2) node[below left, font=\scriptsize, text=black] {$\left(\begin{smallmatrix} 3 \\ -1 \end{smallmatrix}\right)$};
    
    % Correction: Claude Opus 5 (High) --- this caption sat at (2.5,-1.5), which is where the
    % `below left' label of V2 = (3,-1) also lands, so the column vector printed on top of
    % "Area is |det L| = 5". Dropped clear of it; nothing else lies below the parallelogram.
    \node[blue, align=center, font=\small] at (2.5, -2.1) {Area is $\lvert\det L\rvert = 5$\\ times larger};
  \end{scope}
\end{tikzpicture}
\end{center}

The change of variables formula generalizes this to non-linear maps: the Jacobian determinant acts as the local volume scaling factor.
\end{remark}

% Source: Jérôme Paschoud/Notizen/Woche 8.2 Riemann-Integral.pdf, pp. 2--3
% Generator: Gemini 3.1 Pro (High)
\begin{example}[Gaussian integral via polar coordinates]
\label{ex:gaussian_integral_polar}
A classic application of the change of variables formula is computing the Gaussian integral $I := \int_{-\infty}^\infty e^{-x^2}\,dx$. We consider its square:
\[
I^2 = \left( \int_{-\infty}^\infty e^{-x^2}\,dx \right) \left( \int_{-\infty}^\infty e^{-y^2}\,dy \right) = \int_{\mathbb{R}^2} e^{-(x^2+y^2)} \mathop{d\text{vol}}(x,y).
\]
We use polar coordinates, but strictly following \cref{thm:change_of_variables}, we define the forward map $\Phi: \mathbb{R}^2 \setminus \{(0,0)\} \to (0, \infty) \times [-\pi, \pi)$ by
\[
\Phi(x,y) := \left( \sqrt{x^2+y^2}, \operatorname{atan2}(y,x) \right) =: (r, \varphi).
\]

\begin{center}
% Correction: Claude Opus 5 (Max) --- the picture exhibited no size difference at all, under a
% caption reading "not equal in size". The two source rectangles were $[0.5,1.5]\times[0,0.5]$
% and $[0.5,1.5]\times[0.5,1]$: congruent, area $0.5$ each. Their images were the sectors
% $r\in[0.5,1.5]$ over $0$--$30^\circ$ and over $30$--$60^\circ$: also congruent, area $0.5236$
% each. Stacking the two rectangles in $\varphi$ was the cause, because $\det\Jac\Phi^{-1} = r$
% depends on $r$ alone: shifting a rectangle in $\varphi$ cannot change its image's area, and
% shifting it in $r$ is the only thing that can.
%
% The two rectangles are now side by side in $r$, both $0.5 \times 0.6$, at the same
% $\varphi$-band. Their images are then sectors of the same angular width but different radii:
%   red   $\tfrac12(45^\circ-10^\circ)(0.9^2-0.4^2) = \tfrac12(0.6109)(0.65) = 0.199$,
%   blue  $\tfrac12(45^\circ-10^\circ)(2.0^2-1.5^2) = \tfrac12(0.6109)(1.75) = 0.535$,
% a factor of $2.69$, which is the ratio of the mean radii $1.75/0.65$ as $\det\Jac = r$ predicts.
\begin{tikzpicture}[>=Latex, scale=1]
  % Left side: (r, phi) plane. Two CONGRUENT rectangles, 0.5 wide and 0.6 tall, differing only
  % in where they sit along r -- so r_2 - r_1 = r_4 - r_3 = 0.5 by construction.
  \draw[->, gray] (-0.2, 0) -- (2.6, 0) node[right, font=\scriptsize, text=black] {$r$};
  \draw[->, gray] (0, -0.2) -- (0, 1.8) node[above, font=\scriptsize, text=black] {$\varphi$};

  \draw[thick, fill=red!20]  (0.4, 0.5) rectangle (0.9, 1.1);
  \draw[thick, fill=blue!20] (1.5, 0.5) rectangle (2.0, 1.1);

  \node[left, font=\scriptsize] at (0, 1.1) {$\varphi_2$};
  \node[left, font=\scriptsize] at (0, 0.5) {$\varphi_1$};
  \node[below, font=\scriptsize] at (0.4, 0) {$r_1$};
  \node[below, font=\scriptsize] at (0.9, 0) {$r_2$};
  \node[below, font=\scriptsize] at (1.5, 0) {$r_3$};
  \node[below, font=\scriptsize] at (2.0, 0) {$r_4$};
  % Anchored south WEST, not south: centred at x = 1.2 this line is about 2.6 units wide, so it
  % reached back past x = 0 and the $\varphi$ axis printed through its opening word. Rendering
  % p.278 showed it; anchoring the left edge is the only placement that cannot drift back.
  \node[font=\scriptsize, anchor=south west] at (0.15, 1.25) {same $\Delta r$, same $\Delta\varphi$};

  % Arrow
  \draw[->, thick] (3, 0.5) .. controls (3.3, 0.8) and (3.7, 0.8) .. (4, 0.5) node[midway, above] {$\Phi^{-1}$};

  % Right side: (x,y) plane. Same angular band for both, so only the radius differs.
  \begin{scope}[shift={(6.5, 0.5)}]
    \draw[->, gray] (-0.6, 0) -- (2.6, 0) node[right, font=\scriptsize, text=black] {$x$};
    \draw[->, gray] (0, -0.6) -- (0, 2.6) node[above, font=\scriptsize, text=black] {$y$};

    \foreach \r in {0.4, 0.9, 1.5, 2.0} { \draw[gray!45, dashed] (0,0) circle (\r); }

    % Red sector: the image of the rectangle near the origin.
    \draw[thick, fill=red!20]
      (10:0.4) -- (10:0.9) arc (10:45:0.9) -- (45:0.4) arc (45:10:0.4) -- cycle;
    % Blue sector: the image of the congruent rectangle further out.
    \draw[thick, fill=blue!20]
      (10:1.5) -- (10:2.0) arc (10:45:2.0) -- (45:1.5) arc (45:10:1.5) -- cycle;
  \end{scope}

  \node[align=center, font=\small] at (4.4, -1.7)
    {two congruent rectangles, two images of \emph{different} area\\
     $\Rightarrow$ we must compensate with the (Jacobian) determinant};
\end{tikzpicture}
\end{center}

The Jacobian matrix is
\[
\Jac\Phi(x,y) = \begin{pmatrix} \frac{x}{\sqrt{x^2+y^2}} & \frac{y}{\sqrt{x^2+y^2}} \\[1ex] \frac{-y}{x^2+y^2} & \frac{x}{x^2+y^2} \end{pmatrix}.
\]
Its determinant is
\[
\det \Jac\Phi(x,y) = \frac{x^2}{(x^2+y^2)^{3/2}} + \frac{y^2}{(x^2+y^2)^{3/2}} = \frac{1}{\sqrt{x^2+y^2}} = \frac{1}{r}.
\]
Applying the theorem (ignoring the null set $\{(0,0)\}$ where $\Phi$ is undefined), we get
\[
I^2 = \int_{\Phi(\mathbb{R}^2 \setminus \{(0,0)\})} \frac{e^{-r^2}}{\lvert1/r\rvert} \,dr\,d\varphi = \int_{-\pi}^\pi d\varphi \int_0^\infty r e^{-r^2}\,dr = 2\pi \cdot \left[ -\frac{1}{2} e^{-r^2} \right]_0^\infty = 2\pi \cdot \frac{1}{2} = \pi.
\]
Since $I > 0$, we conclude that $I = \sqrt{\pi}$.
\end{example}

% Supplement: Sascha Brack/Class Notes/Week_08_Notes_Friday.pdf, pp. 2--3
\begin{example}[Change of variables on a diamond]
\label{ex:change_of_variables_diamond}
Compute $I := \int_E e^{x+y}(x-y) \mathop{d\text{vol}}(x,y)$, where $E$ is the region bounded by the lines $y=x$, $y=-x+2$, $y=x-2$, and $y=-x$.

Writing out the inequalities, $E = \{ y \leq x \} \cap \{ y \leq -x+2 \} \cap \{ y \geq -x \} \cap \{ y \geq x-2 \}$.
Rearranging these gives $0 \leq x-y \leq 2$ and $0 \leq x+y \leq 2$.
This suggests the substitution $u := x+y$ and $v := x-y$. The new domain is the square $V = [0,2] \times [0,2]$.
The inverse map is $\Phi^{-1}(u,v) = (x,y) = \big(\frac{u+v}{2}, \frac{u-v}{2}\big)$.
Its Jacobian matrix is $\Jac\Phi^{-1}(u,v) = \begin{pmatrix} 1/2 & 1/2 \\ 1/2 & -1/2 \end{pmatrix}$, so $\big|\det\Jac\Phi^{-1}\big| = \big|-\frac{1}{4} - \frac{1}{4}\big| = \frac{1}{2}$.
By the change of variables formula:
\[I = \int_0^2 \int_0^2 e^u v \left|\det\Jac\Phi^{-1}\right| \,du\,dv = \frac{1}{2} \int_0^2 \int_0^2 e^u v \,du\,dv = \frac{1}{2} \left[ e^u \right]_0^2 \left[ \frac{v^2}{2} \right]_0^2 = \frac{1}{2} (e^2-1) \cdot 2 = e^2-1.\]
\end{example}

\begin{example}[Improper integral via polar coordinates]
\label{ex:improper_integral_polar}
Compute $\int_{\mathbb{R}^2} \frac{1}{(1+x^2+y^2)^2} \mathop{d\text{vol}}(x,y)$.
This integral is improper because the domain is unbounded. Since the integrand is non-negative and continuous, its Riemann improper integral is well-defined, and we can compute it using polar coordinates $x = r\cos\theta, y = r\sin\theta$.
The Jacobian determinant is $r$. The domain $\mathbb{R}^2$ corresponds to $r \in [0, \infty)$ and $\theta \in [0, 2\pi)$.
\[
\int_{\mathbb{R}^2} \frac{1}{(1+x^2+y^2)^2} \,dx\,dy
= \int_0^{2\pi} \int_0^\infty \frac{r}{(1+r^2)^2} \,dr\,d\theta.
\]
Substituting $u = r^2$, $du = 2r\,dr$, the inner integral is $\frac{1}{2} \int_0^\infty \frac{1}{(1+u)^2} \,du = \frac{1}{2} \left[ \frac{-1}{1+u} \right]_0^\infty = \frac{1}{2}$.
The outer integral then gives $\frac{1}{2} \int_0^{2\pi} d\theta = \pi$.
\end{example}

% Supplement: Sascha Brack/Class Notes/Week_08_Notes_Monday.pdf, pp. 8--9
\begin{example}[Volume of a sphere]
\label{ex:volume_sphere}
Let $A := \{(x,y,z) \in \mathbb{R}^3 : x^2+y^2+z^2 \leq 1\}$ be the unit ball. Given that $A$ is Jordan measurable, its volume is $\mu(A) = \int_A 1 \mathop{d\text{vol}}(x,y,z)$.
We can compute this using spherical coordinates:
\[\Phi(r,\theta,\varphi) = \begin{pmatrix} r\sin\theta\cos\varphi \\ r\sin\theta\sin\varphi \\ r\cos\theta \end{pmatrix}.\]
Up to a Jordan null set (which does not change the volume), the domain $A$ corresponds to the box $V = (0,1) \times (0,\pi) \times (-\pi,\pi)$.
The Jacobian determinant is $\det\Jac\Phi(r,\theta,\varphi) = r^2\sin\theta$, which is positive on this domain.
Applying the change of variables formula and Fubini's theorem (which allows slicing up the integral):
\begin{align*}
\mu(A) &= \int_{\Phi^{-1}(A)} \big|\det\Jac\Phi(r,\theta,\varphi)\big| \,dr\,d\theta\,d\varphi \\
&= \int_{-\pi}^\pi \int_0^\pi \int_0^1 r^2\sin\theta \,dr\,d\theta\,d\varphi \\
&= \left(\int_{-\pi}^\pi d\varphi\right) \left(\int_0^\pi \sin\theta \,d\theta\right) \left(\int_0^1 r^2 \,dr\right) \\
&= (2\pi) \cdot (2) \cdot \left(\frac{1}{3}\right) = \frac{4\pi}{3}.
\end{align*}
\end{example}

% Quelle: Linus Lüchinger/Slides/Slides-05-11.pdf, slide 3
% Generator: Claude Opus 5 (High)
\begin{remark}[Polar coordinates as slicing into shells]
\label{rem:polar_as_shells}
The factor $r^2$ in the computation above, and the $r$ in the two-dimensional case, are usually
produced by grinding out a Jacobian determinant. There is a reading that gets them without any
computation, and that generalises to $\mathbb{R}^n$ at no extra cost.

Slice $\mathbb{R}^n$ into spheres centred at the origin. The sphere of radius $r$ is an
$(n-1)$-dimensional submanifold, and it is the unit sphere scaled by $r$; since scaling by $r$
multiplies $(n-1)$-dimensional volume by $r^{n-1}$, its volume is
\[\vol_{n-1}\big(r\,\mathbb{S}^{n-1}\big) = C_n\,r^{n-1},
  \qquad C_n := \vol_{n-1}\big(\mathbb{S}^{n-1}\big).\]
A shell between radii $r$ and $r+dr$ is that sphere thickened by $dr$, so it contributes
$C_n r^{n-1}\,dr$ of $n$-dimensional volume. Adding the shells,
\[\int_{B_R(0)} f(\lvert x\rvert)\,dx = C_n\int_0^R f(r)\,r^{n-1}\,dr
  \qquad\text{for radial } f.\]

The powers now need no memorising: $n=2$ gives $r^{1}$ and $n=3$ gives $r^{2}$, matching the two
Jacobians above, and $C_2 = 2\pi$, $C_3 = 4\pi$ are the circumference of the unit circle and the
area of the unit sphere. Taking $f \equiv 1$, $n=3$, $R=1$ reproduces
$4\pi\int_0^1 r^2\,dr = 4\pi/3$ in one line.

This is a genuine shortcut only for \emph{radial} integrands, where the angular integral is a
constant that can be pulled out; when $f$ depends on the angles as well, the full change of
variables is unavoidable. But it explains where the radial power comes from, which the
determinant computation does not.
\end{remark}

% Supplement: Analysis_II_Script_v1.pdf, pp. 113--114
\begin{remark}[Reference: the standard coordinate Jacobians]
\label{rem:coordinate_jacobians_table}
The polar, cylindrical and spherical substitutions recur throughout this chapter and the next
several; collected here in one place, with $r \geq 0$ throughout and each parametrizing its
target domain up to a Jordan null set:
\begin{center}
\begin{tabular}{@{}lll@{}}
\toprule
\textbf{Coordinates} & \textbf{Transformation} & $\lvert\det\Jac\Phi\rvert$ \\
\midrule
Polar ($\mathbb{R}^2$) &
  $\Phi(r,\varphi) = (r\cos\varphi,\ r\sin\varphi)$ & $r$ \\
Cylindrical ($\mathbb{R}^3$) &
  $\Phi(r,\varphi,z) = (r\cos\varphi,\ r\sin\varphi,\ z)$ & $r$ \\
Spherical ($\mathbb{R}^3$) &
  $\Phi(r,\theta,\varphi) = (r\sin\theta\cos\varphi,\ r\sin\theta\sin\varphi,\ r\cos\theta)$
  & $r^2\sin\theta$ \\
\bottomrule
\end{tabular}
\end{center}
The angular ranges are $\varphi \in (-\pi,\pi)$ throughout, and $\theta \in (0,\pi)$
(\newterm{colatitude}, measured from the positive $z$-axis) for the spherical case. Cylindrical
coordinates are just polar coordinates in the first two variables, carried along unchanged in
$z$, which is exactly why the Jacobian is the same $r$ as in the planar case: the extra $dz$
contributes a factor of $1$. Spherical coordinates are used already in
\cref{ex:volume_sphere} above; cylindrical coordinates are the natural choice whenever a region
has rotational symmetry about an axis without being a ball, e.g.\ a solid cylinder or cone.
\end{remark}

% Source: old_exams/examFS24.pdf (21 August 2024), Wahr/Falsch-Fragen, Aufgabe 8, p. 5
% Extractor: Claude Opus 5 (High)
% Worked out here rather than set as an exercise: it is a checklist against the table immediately
% above, and its three-dimensional counterpart ex:examHS24_TF_spherical is already an exercise
% further down this file. Note the exam's argument order (theta, r), which is the reverse of the
% table's; kept as the exam has it, since the sign of the determinant depends on it.
\begin{example}[Polar coordinates, checked item by item]
\label{ex:polar_coordinates_checklist}
% Generator: Claude Opus 5 (High) --- the four statements are the examiner's; the working is
% written for these notes.
Consider $\Psi : \mathbb{R}\times[0,\infty) \to \mathbb{R}^2$,
$\Psi(\theta, r) := (r\cos\theta,\ r\sin\theta)$, and note the argument order: the angle first.
Four statements, of which two are true.
\begin{enumerate}[label=\textbf{(\alph*)}]
  \item \textbf{$\Psi$ restricted to $(0,2\pi]\times(0,\infty)$ is injective: true.} Suppose
    $\Psi(\theta_1,r_1) = \Psi(\theta_2,r_2)$. Taking Euclidean norms gives $r_1 = r_2 =: r > 0$,
    and dividing by $r$ leaves $(\cos\theta_1,\sin\theta_1) = (\cos\theta_2,\sin\theta_2)$, so
    $\theta_1 - \theta_2 \in 2\pi\mathbb{Z}$. A half-open interval of length $2\pi$ contains at
    most one representative of each such class, so $\theta_1 = \theta_2$.
  \item \textbf{$\Psi$ restricted to $(0,2\pi]\times(0,\infty)$ is surjective: false.} Its image is
    $\mathbb{R}^2\setminus\{(0,0)\}$: every non-zero point is reached, by \textbf{(a)} exactly
    once, and the origin is reached only from $r = 0$, which the restriction excludes.
  \item \textbf{$\det\Jac\Psi(\theta,r) = r^2\sin\theta\cos\theta$: false.} The columns of the
    Jacobi matrix are the partial derivatives in $\theta$ and in $r$, in that order:
    \[\Jac\Psi(\theta,r) = \begin{pmatrix} -r\sin\theta & \cos\theta \\
      r\cos\theta & \sin\theta\end{pmatrix}, \qquad
      \det\Jac\Psi(\theta,r) = -r\sin^2\theta - r\cos^2\theta = -r .\]
    The magnitude $\lvert\det\Jac\Psi\rvert = r$ is the entry in
    \cref{rem:coordinate_jacobians_table}; the minus sign appears only because the exam lists the
    angle before the radius, which swaps two columns.
  \item \textbf{$\vol_2(B_R) = 2\pi\int_0^R r\,dr$ for $B_R = \{x^2+y^2 < R^2\}$: true.} By
    \cref{thm:change_of_variables} with the weight $\lvert\det\Jac\Psi\rvert = r$, and since the
    integrand $1$ does not depend on $\theta$, the angular integral contributes a factor $2\pi$.
    The right-hand side is $2\pi\cdot R^2/2 = \pi R^2$, as it should be. This is also the $n=2$
    case of \cref{rem:polar_as_shells}.
\end{enumerate}
The proposed formula in \textbf{(c)} is worth a second look, because it is not a random wrong
answer: $r^2\sin\theta$ is the \emph{spherical} Jacobian, and $\cos\theta$ is what a misremembered
product rule adds to it. There is a cheap guard against importing one row of the table into
another. Every system in it has the form $\Phi(r,\omega) = r\,u(\omega)$ for a map $u$ into the
unit sphere, so $\Phi$ is homogeneous of degree $1$ in $r$. The radial column of $\Jac\Phi$ is then
$u(\omega)$, free of $r$, while each of the $n-1$ angular columns is $r\,\partial u/\partial\omega_i$
and so carries exactly one factor $r$. Multilinearity of the determinant pulls all $n-1$ of them
out, leaving $\lvert\det\Jac\Phi\rvert = r^{n-1}$ times a function of the angles alone. In
$\mathbb{R}^2$ that power is $r^1$, and the proposal carries $r^2$.
\exinfo{This example is taken from the exam of 21 August 2024 (Prof.\ Joaquim Serra), where it appears as the eighth block
of true/false questions.}
\end{example}

% Source: Jérôme Paschoud/Notizen/Woche 9.1 Variablenwechsel.pdf, p. 2
% Generator: Gemini 3.1 Pro (High)
\begin{example}[Volume of an ellipse]
\label{ex:volume_ellipse}
Compute the area (2-dimensional volume) of the ellipse $E := \{(x,y) \in \mathbb{R}^2 \mid (x/a)^2 + (y/b)^2 < 1\}$ for $a,b > 0$.
We use the transformation $\Phi(r,\varphi) = (ar\cos\varphi, br\sin\varphi)$. This maps the rectangle $V = (0,1) \times (0,2\pi)$ bijectively onto $E$ minus a null set (the positive $x$-axis).
The Jacobian matrix is
\[\Jac\Phi(r,\varphi) = \begin{pmatrix} a\cos\varphi & -ar\sin\varphi \\ b\sin\varphi & br\cos\varphi \end{pmatrix},\]
and its determinant is $\det\Jac\Phi(r,\varphi) = abr(\cos^2\varphi + \sin^2\varphi) = abr$.
Thus, by the change of variables formula,
\[\int_E 1 \,dx\,dy = \int_0^{2\pi} \int_0^1 abr \,dr\,d\varphi = \left(\int_0^{2\pi} ab \,d\varphi\right) \left(\int_0^1 r \,dr\right) = 2\pi ab \cdot \frac{1}{2} = \pi ab.\]
\end{example}

% Source: Jérôme Paschoud/Notizen/Woche 9.1 Variablenwechsel.pdf, p. 2
% Generator: Gemini 3.1 Pro (High)
\begin{example}[Area of a cardioid]
\label{ex:area_cardioid}
Compute the area of the region $\Omega$ given in polar coordinates by $\Omega := \{(r,\varphi) \mid \varphi \in [0, 2\pi),\ 0 \leq r \leq 1 + \cos\varphi\}$.
Using polar coordinates $\Phi(r,\varphi) = (r\cos\varphi, r\sin\varphi)$, the area is
\[\int_\Omega 1 \,dx\,dy = \int_0^{2\pi} \int_0^{1+\cos\varphi} r \,dr\,d\varphi = \int_0^{2\pi} \left[ \frac{r^2}{2} \right]_{r=0}^{r=1+\cos\varphi} \,d\varphi = \frac{1}{2} \int_0^{2\pi} (1+\cos\varphi)^2 \,d\varphi.\]
Expanding the integrand: $(1+\cos\varphi)^2 = 1 + 2\cos\varphi + \cos^2\varphi$. We know $\int_0^{2\pi} \cos\varphi \,d\varphi = 0$ and $\int_0^{2\pi} \cos^2\varphi \,d\varphi = \pi$. Thus,
\[\int_\Omega 1 \,dx\,dy = \frac{1}{2} (2\pi + 0 + \pi) = \frac{3\pi}{2}.\]
\end{example}

% Source: Jérôme Paschoud/Notizen/Woche 9.1 Variablenwechsel.pdf, p. 3
% Generator: Gemini 3.1 Pro (High)
\begin{example}[Combining substitution and Fubini]
\label{ex:substitution_and_fubini_rational}
Compute $\int_\Omega (x^3 y + x y^3) \,dx\,dy$ where $\Omega := \{(x,y) \in \mathbb{R}^2 \mid x,y > 0,\ 1/x \le y \le 3/x,\ 1 \le x^2 - y^2 \le 4\}$.
The integrand factors as $xy(x^2+y^2)$. The boundary inequalities are equivalent to $1 \le xy \le 3$ and $1 \le x^2-y^2 \le 4$, which strongly suggests defining $u := xy$ and $v := x^2 - y^2$.
The domain in the $(u,v)$-plane is the rectangle $R = [1,3] \times [1,4]$.
Let $\Psi(x,y) = (xy, x^2-y^2) = (u,v)$. The Jacobian of this inverse transformation is
\[\Jac\Psi(x,y) = \begin{pmatrix} y & x \\ 2x & -2y \end{pmatrix}, \qquad \det\Jac\Psi(x,y) = -2y^2 - 2x^2 = -2(x^2+y^2).\]
For the forward transformation $\Phi = \Psi^{-1}$, we have $\big|\det\Jac\Phi(u,v)\big| = \frac{1}{\lvert\det\Jac\Psi(x,y)\rvert} = \frac{1}{2(x^2+y^2)}$.
Applying the change of variables formula:
\[\int_\Omega xy(x^2+y^2) \,dx\,dy = \int_R u(x^2+y^2) \frac{1}{2(x^2+y^2)} \,du\,dv = \frac{1}{2} \int_1^4 \int_1^3 u \,du\,dv = \frac{1}{2} \cdot 3 \cdot \left[ \frac{u^2}{2} \right]_1^3 = \frac{3}{2} \cdot 4 = 6.\]
\end{example}


% Source: Corsin Nick/Class Notes/Week 8.pdf, p. 7


% --- exercises relocated here from the dissolved sheet 8 ---

% Originally: content/week-08.tex, \exercisesheet{8}, problem 8.5
% Source: exercises/Ex8_Analysis2_eng.pdf

% Source: exercises/Ex8_Analysis2_eng.pdf, pp. 1--2
\begin{exercise}[Change of variables and Jacobians \textnormal{(important)}]
\label{ex:change_of_variables_computations}
For each of the following domains and change of variables find the jacobian and the appropriate
transformed domain. There is no need to actually compute the integrals!
\begin{enumerate}[label=\textbf{(\alph*)}]
  \item $A := \{(x_1,x_2) \in \mathbb{R}^2 : x_1 > 0,\ x_2 < x_1,\ 1 < x_1^2+x_2^2 < 4\}$ and
    $x_1 = r\cos\theta$, $x_2 = r\sin\theta$. Using the change of variables formula, complete
    the dots in the following formula
    \[\int_A x_1^2\sin(x_2)\,dx_1dx_2 = \int_{\dots}\dots\,dr\,d\theta\]
  \item $B := \{(x,y) \in \mathbb{R}^2 \mid 1 < xy < 2,\ x^2 < y < 2x^2\}$ and $u := xy$,
    $v := x^2$. Using the change of variables formula, complete the dots in the following
    formula
    \[\int_B y^2e^{-xy}\,dx\,dy = \int_{\dots}\dots\,du\,dv\]
  \item $C := \{(x,y,z) \in \mathbb{R}^3 \mid 1 < z-2y < 2,\ 0 < z < 1,\ -2 < 3x+y+z < 0\}$ and
    $u := z$, $v := z-2y$, $w := 3x+y+z$. Using the change of variables formula, complete the
    dots in the following formula
    \[\int_C xyz\,dx\,dy = \int_{\dots}\dots\,du\,dv\,dw.\]
\end{enumerate}

% Supplement: Sascha Brack/Ex Sheet Hints/Ex8_Analysis2_hints.pdf, p. 1
\exhint[Sascha Brack's hint on 8.5.2]{First pin down the signs. The constraints force
$x \neq 0$ and $y \neq 0$; then $y > 0$ forces $x > 0$, so $u > 0$ and $v > 0$ throughout $B$.
That is what lets you divide by $x$ and $y$ and rearrange the inequalities without their
directions flipping.}
\exinfo{This exercise is Problem 8.5 of Problem Sheet 8, Analysis II, Spring Semester 2026.}
\exsol{sol:change_of_variables_computations}
\end{exercise}

\begin{ainote}
The integral in \textbf{8.5.3} is written $\int_C xyz\,dx\,dy$ on the official sheet, but $C$ is
a subset of $\mathbb{R}^3$ and the right-hand side carries three differentials, so this must
read $dx\,dy\,dz$. Quoted verbatim above, as the sheet is a primary source.
\end{ainote}

% An AI-Note stood here until 2026-08-09 announcing, above the exercise, that part (c)'s domain is
% the box (0,1)x(1,2)x(-2,0) and that part (b)'s second constraint becomes v^{3/2} < u < 2v^{3/2}.
% That is the answer to two of the three parts. Moved into 99-solutions.tex per the spoiler rule
% in style.md; the surviving half duplicated the AI-Note already sitting beside the solution.

\begin{aiexercise}[A substitution read off the integrand]
% Generator: Claude Opus 5 (High)
\label{ex:ai_triangle_substitution}
Let $T$ be the closed triangle with vertices $(0,0)$, $(1,0)$ and $(0,1)$, and consider
\[I := \int_T \exp\!\left(\frac{y-x}{y+x}\right)dx\,dy.\]
\begin{enumerate}[label=\textbf{(\alph*)}]
  \item Substitute $u := y-x$ and $v := y+x$. Show that the interior of $T$ corresponds to
    $\{(u,v) : 0 < v < 1,\ -v < u < v\}$, and find the factor relating $dx\,dy$ to $du\,dv$.
  \item Evaluate $I$. The inner integral is the one to do first, and it is elementary.
  \item The integrand is undefined at the origin, which lies in $T$. Why does this not affect
    the answer?
\end{enumerate}
\exsol{sol:ai_triangle_substitution}
\end{aiexercise}

% Source: old_exams/examHS24.pdf (13 February 2025), Kurzproblem 1, part 1, p. 8
% Extractor: Gemini 3.6 Flash (High)
% Correction: Claude Opus 5 (High) --- \operatorname{vol} -> \vol and J\Phi -> \Jac\Phi per the
% declared-macro rule.
\begin{exercise}[Volume of an ellipsoid via linear change of variables]
\label{ex:volume_ellipsoid_change_variables}
Let $a, b, c > 0$ and consider the solid ellipsoid
\[U := \left\{(x,y,z) \in \mathbb{R}^3 \ \middle|\ \left(\frac{x}{a}\right)^2 + \left(\frac{y}{b}\right)^2 + \left(\frac{z}{c}\right)^2 \leq 1\right\}.\]
\begin{enumerate}[label=\textbf{(\alph*)}]
  \item Define a linear change of variables $\Phi : \mathbb{R}^3 \to \mathbb{R}^3$ mapping the closed unit ball $\overline{B_1(0)} \subseteq \mathbb{R}^3$ bijectively onto $U$.
  \item Compute the Jacobian determinant $\det \Jac\Phi(u,v,w)$ of this transformation.
  \item Compute the three-dimensional Jordan volume $\vol_3(U)$ using the change of variables formula and the known volume of the unit ball.
\end{enumerate}
\exinfo{This exercise is taken from the exam of 13 February 2025 (Prof.\ Joaquim Serra), Short Problem 1, part 1, which
also asks for a labelled 3D sketch of $U$.}
\exsol{sol:volume_ellipsoid_change_variables}
\end{exercise}

\begin{ainote}
The original exam writes the ellipsoid with an equality, $(x/a)^2+(y/b)^2+(z/c)^2 = 1$, and then
asks for its volume --- but that set is the ellipsoid \emph{surface}, whose Jordan volume is $0$.
The intended object is the solid region it bounds, so the inequality is used here.
\end{ainote}

% Source: old_exams/fs2023/Prüfung.pdf (9 August 2021), Teil A, Frage 5, p. 6
% Extractor: Claude Opus 5 (High)
\begin{exercise}[A cubic shear, and the area it sweeps out]
\label{ex:cubic_shear_area}
Let $\varphi : \mathbb{R}^2 \to \mathbb{R}^2$ be given by $\varphi(x) := (x_1^3 - x_2,\; x_2 + x_1)$,
where $x = (x_1,x_2)$.
\begin{enumerate}[label=\textbf{(\alph*)}]
  \item Compute the derivative $D_x\varphi$ for every $x \in \mathbb{R}^2$.
  \item Show that $\varphi : \mathbb{R}^2 \to \varphi(\mathbb{R}^2)$ is a diffeomorphism.
  \item What is the area of the image $\varphi(Q)$ of the square $Q := [0,1] \times [0,1]$?
\end{enumerate}
\exinfo{This exercise is taken from the Analysis I \& II exam of 9 August 2021 (Prof.\ Giovanni Felder), Part A, Problem 5.}
\exsol{sol:cubic_shear_area}
\end{exercise}

% Source: old_exams/examHS24.pdf (13 February 2025), Wahr/Falsch-Fragen, Aufgabe 9, p. 5
% Extractor: Claude Opus 5 (High)
\begin{exercise}[A Gaussian second moment in polar coordinates]
\label{ex:gaussian_second_moment_polar}
For $\alpha > 0$, compute the integral
\[I_\alpha := \int_{\mathbb{R}^2} x^2 e^{-\alpha(x^2+y^2)} \, dx\,dy .\]
\exhint[Official hint]{Notice that, by symmetry, $I_\alpha = \int_{\mathbb{R}^2} y^2
e^{-\alpha(x^2+y^2)}\,dx\,dy$ as well. Compute $2I_\alpha$ using polar coordinates.}
\exinfo{This exercise is taken from the exam of 13 February 2025 (Prof.\ Joaquim Serra), where it appears as the ninth
block of true/false questions, there phrased as a series of checkpoints along this computation.}
\exsol{sol:gaussian_second_moment_polar}
\end{exercise}

% Source: old_exams/examFS24.pdf (21 August 2024), Wahr/Falsch-Fragen, Aufgabe 9, p. 5
% Extractor: Claude Opus 5 (High)
% Filed next to its February 2025 sibling above deliberately: both are one integral over R^2 in
% polar coordinates, and the pair is worth doing together because the weight differs (Gaussian
% against exponential-of-the-radius) while the method does not.
\begin{exercise}[An exponential of the radius, and how it decays in $\alpha$]
\label{ex:radial_exponential_integral}
For $\alpha > 0$, compute
\[I_\alpha := \int_{\mathbb{R}^2} e^{-\alpha\sqrt{x^2+y^2}}\,dx\,dy,\]
and then decide whether each of the following is true or false: $I_1 = \pi$; $I_2 = \pi/2$;
$\lim_{\alpha\to+\infty}\alpha I_\alpha = 0$.
\exhint[Official hint]{You may use polar coordinates and integration by parts.}
\exinfo{This exercise is taken from the exam of 21 August 2024 (Prof.\ Joaquim Serra), where it appears as the ninth block
of true/false questions. The instruction to compute $I_\alpha$ is the exam's own preamble to those
three statements.}
\exsol{sol:radial_exponential_integral}
\end{exercise}

% Source: old_exams/examHS24.pdf (13 February 2025), Wahr/Falsch-Fragen, Aufgabe 6, p. 4
% Extractor: Claude Opus 5 (High)
% Filed in this chapter rather than in 15-inverse-function-theorem, where a Gemini 3.1 Pro pass had
% put a different (and invented) block under this citation: only part (a) below is the inverse
% function theorem, and parts (b) and (c) are change of variables.
\begin{exercise}[True or False \textnormal{(the image of a small set under a $C^1$ map)}]
\label{ex:examHS24_TF_image_of_small_sets}
Consider $\Phi \in C^1(\mathbb{R}^2, \mathbb{R}^2)$ with
\[\Phi(0,0) = \begin{pmatrix} 0 \\ 0 \end{pmatrix} \quad\text{and}\quad
  \Jac\Phi(0,0) = \begin{pmatrix} 2 & 1 \\ -1 & 1 \end{pmatrix},\]
and let $X := \Phi^{-1}\big((0,0)\transp\big)$. Decide whether each of the following statements is
true or false.
\begin{enumerate}[label=\textbf{(\alph*)}]
  \item For $r > 0$ small enough, the set $\Phi(B_r(0,0))$ is open.
  \item $\vol_2\big(\Phi(B_r(0,0))\big) = 2\pi r^2 + o(r^2)$ as $r \to 0^+$.
  \item $\vol_2\big(\Phi([0,r] \times [0, ar^{1/2}])\big) = 3ar^{3/2} + o(r^{3/2})$ as
    $r \to 0^+$, where $a > 0$ is a constant.
\end{enumerate}
\exinfo{This exercise is taken from the exam of 13 February 2025 (Prof.\ Joaquim Serra), where it appears as the sixth
block of true/false questions.}
\exsol{sol:image_of_small_sets_tf}
\end{exercise}

% Source: old_exams/examHS24.pdf (13 February 2025), Wahr/Falsch-Fragen, Aufgabe 7, p. 4
% Extractor: Claude Opus 5 (High)
% Replaces a block that a Gemini 3.1 Pro pass filed under this same citation. That block asked for
% the volume of an ellipsoid spanned by three vectors of determinant 2, which is not on p. 4 and is
% nowhere on this exam; it survives, honestly relabelled, as ex:ai_spanned_ellipsoid_volume below.
\begin{exercise}[True or False \textnormal{(volumes under an affine map)}]
\label{ex:examHS24_TF_affine_volume_scaling}
Consider the map
\[\Psi : \mathbb{R}^3 \to \mathbb{R}^3, \qquad
  \Psi(x,y,z) := \big(-x + e^\pi,\; y + \alpha z,\; \alpha y + z\big),\]
where $\alpha \in \mathbb{R}$ is a parameter. Decide whether each of the following statements is
true or false.
\begin{enumerate}[label=\textbf{(\alph*)}]
  \item $\vol_3(\Psi(E)) = 0$ for all Jordan measurable $E \subseteq \mathbb{R}^3$ if and only if
    $\alpha = 1$.
  \item For $\alpha = -2$, $\vol_3(\Psi^k(E)) = 3^k \vol_3(E)$ for all Jordan measurable
    $E \subseteq \mathbb{R}^3$ and all $k \ge 0$, where $\Psi^0 := \id$ and
    $\Psi^{k+1} := \Psi \circ \Psi^k$.
  \item For $\alpha = 0$, $\vol_3(\Psi^{2025}(E)) = \vol_3(E)$ for all Jordan measurable
    $E \subseteq \mathbb{R}^3$.
  \item For all $\alpha \in (0,1)$, $\vol_3\big(\Psi^{-1}([0,1]^3)\big) = -1/(1-\alpha^2)$.
\end{enumerate}
\exinfo{This exercise is taken from the exam of 13 February 2025 (Prof.\ Joaquim Serra), where it appears as the seventh
block of true/false questions.}
\exsol{sol:affine_volume_scaling_tf}
\end{exercise}

% Source: old_exams/examFS24.pdf (21 August 2024), Wahr/Falsch-Fragen, Aufgabe 7, pp. 4--5
% Extractor: Claude Opus 5 (High)
% Worked out rather than set as an exercise, because the exercise directly above is the same
% question with a parameter in it and is strictly harder. What this block adds is the composition
% and inverse rules read off one determinant, which the February 2025 version states only for
% iterates of a single map.
\begin{example}[One determinant, three volume statements]
\label{ex:affine_volume_one_determinant}
% Generator: Claude Opus 5 (High) --- the three statements are the examiner's; the working is
% written for these notes.
Consider $\Psi : \mathbb{R}^3 \to \mathbb{R}^3$, $\Psi(x,y,z) := (z-1,\ x+2y,\ x-y)$. Its linear
part has matrix
\[A = \begin{pmatrix} 0 & 0 & 1 \\ 1 & 2 & 0 \\ 1 & -1 & 0\end{pmatrix},
  \qquad \det A = 1\cdot\det\begin{pmatrix} 1 & 2 \\ 1 & -1\end{pmatrix} = -3,\]
expanding along the first row. The translation by $(-1,0,0)$ does not affect volumes, so
$\vol_3(\Psi(E)) = 3\vol_3(E)$ for every Jordan measurable $E$, by
\cref{thm:change_of_variables}. Everything else follows from the factor $3$.
\begin{enumerate}[label=\textbf{(\alph*)}]
  \item \textbf{$\vol_3(\Psi^{-1}(E)) = \vol_3(E)$ for all Jordan measurable $E$: false.} The
    inverse scales by $\lvert\det A^{-1}\rvert = 1/3$, so $\vol_3(\Psi^{-1}(E)) = \vol_3(E)/3$.
  \item \textbf{$\vol_3(\Psi(E)) = 2\vol_3(E)$: false.} The factor is $3$, not $2$.
  \item \textbf{$\vol_3(\Psi(\Psi(E))) = 9\vol_3(E)$: true.} Applying \textbf{(b)}'s correct factor
    twice gives $3 \cdot 3 = 9$, which is $\lvert\det A^2\rvert = \lvert\det A\rvert^2$.
\end{enumerate}
Only \textbf{(b)} requires the determinant to be computed. Items \textbf{(a)} and \textbf{(c)} then
follow from multiplicativity of the determinant alone, without recomputing anything, which is the
point of the block.
\exhint[Official hint]{To answer these three, it can be useful to compute the Jacobian determinant
of $\Psi$ first.}
\exinfo{This example is taken from the exam of 21 August 2024 (Prof.\ Joaquim Serra), where it appears as the seventh
block of true/false questions.}
\end{example}

% Source: old_exams/examHS24.pdf (13 February 2025), Wahr/Falsch-Fragen, Aufgabe 8, pp. 4--5
% Extractor: Gemini 3.1 Pro (High)
\begin{exercise}[True or False \textnormal{(spherical coordinates)}]
\label{ex:examHS24_TF_spherical}
Consider the polar (spherical) coordinates $\Psi \colon (0, 2\pi) \times (0, \pi) \times (0, \infty)
\to \mathbb{R}^3$ defined by
\[ \Psi(\phi, \theta, r) = (r \cos \phi \sin \theta, r \sin \phi \sin \theta, r \cos \theta). \]
Decide whether each of the following statements is true or false.
\begin{enumerate}[label=\textbf{(\alph*)}]
  \item $\Psi$ is injective.
  \item $\Psi$ is surjective.
  \item The Jacobian determinant $\det \Jac\Psi(\theta, r)$ is given by $r^2 \sin \theta$.
  \item $\vol_3(B_R \setminus B_r) = \pi(R^3 - r^3)$.
  \item $\vol_3(B_R) = 2\pi \int_0^\pi \int_0^R r^2 \sin \theta \, dr \, d\theta$.
  \item $\lim_{R \to 1^+} \frac{\vol_3(B_R \setminus B_1)}{R - 1} = 4\pi$.
\end{enumerate}
\exinfo{This exercise is taken from the exam of 13 February 2025 (Prof.\ Joaquim Serra), where it appears as the eighth
block of true/false questions.}
\exsol{sol:spherical_coordinates_integral_tf}
\end{exercise}

% Source: old_exams/examHS24.pdf (13 February 2025), Kästchenaufgabe 2, p. 7
% Extractor: Gemini 3.1 Pro (High)
\begin{exercise}[Box question: Linear map scaling volume]
\label{ex:examHS24_box_linear_vol}
Give an explicit example of a linear map $\Phi \colon \mathbb{R}^2 \to \mathbb{R}^2$ such that
\[ \vol_2(\Phi(E)) = 2 \vol_2(E) \quad \text{for all } E \subset \mathbb{R}^2 \text{ Jordan measurable.} \]
\exinfo{This exercise is taken from the exam of 13 February 2025 (Prof.\ Joaquim Serra), Box Problem 2, where only the
final answer was graded.}
\exsol{sol:linear_map_volume_scaling_box}
\end{exercise}

% Source: old_exams/examFS24.pdf (21 August 2024), Kästchenaufgabe 2, p. 7
% Extractor: Claude Opus 5 (High)
% The August 2024 counterpart of the box problem directly above: same one-line format, factor 1
% instead of factor 2, and the added clause "different from the identity", which is what stops the
% obvious answer and makes the problem worth having next to its sibling.
\begin{exercise}[Box question: a linear map that preserves every volume]
\label{ex:box_linear_volume_preserving}
Give an explicit example of a linear map $\Phi : \mathbb{R}^2 \to \mathbb{R}^2$, \textbf{different
from the identity}, such that
\[\vol_2(\Phi(E)) = \vol_2(E) \quad \text{for all Jordan measurable } E \subseteq \mathbb{R}^2 .\]
\exinfo{This exercise is taken from the exam of 21 August 2024 (Prof.\ Joaquim Serra), Box Problem 2, where only the final
answer was graded.}
\exsol{sol:box_linear_volume_preserving}
\end{exercise}

% Correction: Claude Opus 5 (High) --- provenance repair. A Gemini 3.1 Pro pass filed the statement
% below as an exam problem, citing old_exams/examHS24.pdf, Aufgabe 7, p. 4. That page carries
% ex:examHS24_TF_affine_volume_scaling instead, and nothing resembling this problem is anywhere on
% that exam. The mathematics is sound, so the block is kept as what it actually is: an exercise
% authored for these notes. Its % Source:, % Extractor: and \exinfo lines have been removed.
\begin{aiexercise}[The volume of a spanned ellipsoid]
\label{ex:ai_spanned_ellipsoid_volume}
% Generator: Gemini 3.1 Pro (High)
Let $v_1, v_2, v_3 \in \mathbb{R}^3$ be vectors with $\det(v_1, v_2, v_3) = 2$, and consider the
region
\[ A := \big\{ t_1 v_1 + t_2 v_2 + t_3 v_3 \ \big|\ (t_1, t_2, t_3) \in \mathbb{R}^3
  \text{ with } t_1^2 + t_2^2 + t_3^2 \le 1 \big\}. \]
Decide whether each of the following statements is true or false.
\begin{enumerate}[label=\textbf{(\alph*)}]
  \item $\vol_3(A) = \frac{8\pi}{3}$.
  \item $A$ is Jordan measurable.
  \item The volume of $A$ depends on the choice of the vectors $v_1, v_2, v_3$.
\end{enumerate}
\exsol{sol:ai_spanned_ellipsoid_volume}
\end{aiexercise}

% Source: old_exams/HS18.pdf (Analysis I & II, winter 2018/2019, Prof. Manfred Einsiedler),
% Aufgabe 13, p. 4
% Extractor: Claude Opus 5 (High)
% This is the proof of thm:jordan_measure_linear_transformation, which chapter 19 states without
% one. Filed here rather than beside the theorem because the exam frames it as the linear case of
% the substitution rule, which is this chapter's subject.
\begin{exercise}[The substitution rule in the linear case]
\label{ex:linear_substitution_elementary_matrices}
Let $A$ be an invertible $n\times n$ matrix. \Cref{thm:jordan_measure_linear_transformation}
asserts that for every Jordan measurable $B \subseteq \mathbb{R}^n$ the image $A(B)$ is again
Jordan measurable, with
\begin{equation}
\label{eq:linear_substitution_rule}
\vol\big(A(B)\big) = \lvert\det A\rvert \, \vol(B).
\end{equation}
This exercise proves it.
\begin{enumerate}[label=\textbf{(\alph*)}]
  \item Prove \eqref{eq:linear_substitution_rule} directly in the case that $A$ is a diagonal
    matrix.
  \item Explain in a few sentences why it now suffices to prove \eqref{eq:linear_substitution_rule}
    for permutation matrices and for shear matrices.
  \item Prove \eqref{eq:linear_substitution_rule} when $A$ is a shear matrix. You may assume for
    simplicity that $B = Q$ is a closed axis-parallel box.
\end{enumerate}
\exinfo{This exercise is taken from the Analysis I \& II examination of winter 2018/2019
(Prof.\ Manfred Einsiedler), Problem 13.}
\exsol{sol:linear_substitution_elementary_matrices}
\end{exercise}

% Source: old_exams/HS17.pdf (Analysis I & II, winter 2017/2018, Prof. Manfred Einsiedler),
% Aufgabe 13, pp. 3--4
% Extractor: Claude Opus 5 (High)
% Kept as a worked example, not an exercise: this is the one genuinely hard lemma behind
% thm:change_of_variables, which this document states without proof, and it is well beyond what
% any exercise here asks. The exam splits it into four graded parts; those are preserved as the
% four steps below.
% Notation normalised to the house convention: the exam writes ||.||_2 for the Euclidean norm,
% which is single bars here (style.md, "Norms"). ||.||_infty keeps its double bars.
\begin{example}[The covering lemma behind the change of variables formula]
\label{ex:cube_covering_lemma_substitution}
\Cref{thm:change_of_variables} is stated in this chapter without proof. The step that makes the
proof hard is not the Jacobian factor but a prior, purely geometric one: a diffeomorphism that is
close to the identity cannot let a cube shrink away. Made precise, that is the following.

\textbf{Claim.} Let $\ell > 0$, let $Q_0 := [-\ell,\ell]^n \subseteq \mathbb{R}^n$, let
$X \supseteq Q_0$ be open and let $\Phi : X \to \Phi(X) \subseteq \mathbb{R}^n$ be a
diffeomorphism with $\Phi(0) = 0$ and $D_0\Phi = I_n$. Suppose $\sigma > 0$ is such that
\[\big\lvert (D_x\Phi - I_n)v \big\rvert \le \sigma \lvert v\rvert
\qquad \text{for all } x \in Q_0,\ v \in \mathbb{R}^n,\]
and assume $s := \sigma\sqrt{n} < 1$. Then
\[(1-s)\,Q_0 \subseteq \Phi(Q_0).\]

Fix $y \in (1-s)Q_0$ and define $F_y : Q_0 \to \mathbb{R}^n$ by
$F_y(x) := y - \big(\Phi(x) - x\big)$.

\begin{enumerate}[label=\textbf{(\alph*)}]
  \item \textbf{Why a fixed point is what we want.} By construction $F_y(x) = x$ holds if and only
    if $y - \Phi(x) + x = x$, that is, if and only if $\Phi(x) = y$. So a fixed point of $F_y$ in
    $Q_0$ is exactly a preimage of $y$ under $\Phi$ lying in $Q_0$, which is what
    $y \in \Phi(Q_0)$ means.

  \item \textbf{The displacement $\Phi - \id$ is a contraction.} Write $\Psi(x) := \Phi(x) - x$, so
    that $D_x\Psi = D_x\Phi - I_n$ and the hypothesis reads
    $\lvert D_x\Psi(v)\rvert \le \sigma\lvert v\rvert$ on $Q_0$, that is,
    $\lVert D_x\Psi\rVert_{\mathrm{op}} \le \sigma$ there. Since $Q_0$ is convex, the segment from
    $x_1$ to $x_2$ stays inside it, and the mean value inequality
    (\cref{thm:mean_value_inequality}) gives
    \[\big\lvert \Phi(x_2)-x_2 - \big(\Phi(x_1)-x_1\big)\big\rvert
      = \lvert \Psi(x_2)-\Psi(x_1)\rvert \le \sigma \lvert x_2-x_1\rvert\]
    for all $x_1,x_2 \in Q_0$.

  \item \textbf{$F_y$ maps $Q_0$ into itself.} Applying \textbf{(b)} with $x_1 := 0$, and using
    $\Psi(0) = \Phi(0)-0 = 0$, we get $\lvert \Psi(x)\rvert \le \sigma\lvert x\rvert$ for
    $x \in Q_0$. For $x \in Q_0$ we have $\lVert x\rVert_\infty \le \ell$ and hence
    $\lvert x\rvert \le \sqrt{n}\,\ell$, so
    \[\lVert \Psi(x)\rVert_\infty \le \lvert \Psi(x)\rvert \le \sigma\sqrt{n}\,\ell = s\ell.\]
    Since $y \in (1-s)Q_0$ means $\lVert y\rVert_\infty \le (1-s)\ell$, the triangle inequality
    gives
    \[\lVert F_y(x)\rVert_\infty \le \lVert y\rVert_\infty + \lVert \Psi(x)\rVert_\infty
      \le (1-s)\ell + s\ell = \ell,\]
    that is, $F_y(x) \in Q_0$. So $F_y : Q_0 \to Q_0$.

  \item \textbf{Conclusion.} $Q_0$ is a closed subset of $\mathbb{R}^n$, hence complete, and by
    \textbf{(b)} the map $F_y$ satisfies
    $\lvert F_y(x_2)-F_y(x_1)\rvert = \lvert \Psi(x_2)-\Psi(x_1)\rvert \le
    \sigma\lvert x_2-x_1\rvert$ with $\sigma \le s < 1$. It is therefore a contraction of the
    complete metric space $Q_0$ into itself, and the Banach fixed point theorem
    (\cref{thm:banach_fixed_point}) supplies a fixed point $x \in Q_0$. By \textbf{(a)},
    $\Phi(x) = y$, so $y \in \Phi(Q_0)$. As $y \in (1-s)Q_0$ was arbitrary,
    $(1-s)Q_0 \subseteq \Phi(Q_0)$.
\end{enumerate}

\begin{remark}[What the lemma is doing in the proof of the change of variables formula]
The formula is proved by cutting the domain into tiny cubes and comparing $\vol(\Phi(Q))$ with
$\lvert\det D\Phi\rvert\vol(Q)$ on each. The upper bound is the easy half: $\Phi$ is Lipschitz on a
compact piece, so $\Phi(Q)$ cannot be much bigger than $Q$. The lower bound is where a naive
argument stalls, because nothing about a Lipschitz map prevents an image from collapsing. This
lemma is that lower bound. First normalise so that $D_0\Phi = I_n$, which is always possible by
composing with a linear map, the case already settled in
\cref{ex:linear_substitution_elementary_matrices}. The lemma then says the image of a cube still
contains a cube shrunk only by the factor $1-s$, and $s \to 0$ as the cubes shrink, because
$D\Phi$ is continuous. Letting the two bounds close is the proof.

Notice which two theorems do the work: the mean value inequality on a convex set in \textbf{(b)},
and the Banach fixed point theorem in \textbf{(d)}. The same pair proves the inverse function
theorem (\cref{thm:inverse_function_theorem}), and that is not a coincidence. This claim is a
quantitative version of the statement that $\Phi$ is locally surjective near a point where its
differential is invertible.
\end{remark}
\exinfo{This example is taken from the Analysis I \& II examination of winter 2017/2018
(Prof.\ Manfred Einsiedler), Problem 13, where the four steps above are set as four graded parts.
The framing and the closing remark are added here.}
\end{example}


% Originally: content/week-08.tex, \subsection{Fubini's theorem} --- promoted to a section

\section{Fubini's theorem}
\label{sec:fubini}

\begin{theorem}[Fubini]
\label{thm:fubini}
Let $X \subseteq \mathbb{R}^m$ and $Y \subseteq \mathbb{R}^n$ be intervals (\qt{boxes}), and let
$f : X\times Y \to \mathbb{R}$, $(x,y)\mapsto f(x,y)$, be continuous. Then
\[\int_{X\times Y} f(x,y)\,dx\,dy
  = \int_X\left(\int_Y f(x,y)\,dy\right)dx
  = \int_Y\left(\int_X f(x,y)\,dx\right)dy,\]
where we look at the expression in brackets as a function, e.g.\
$x \mapsto \int_Y f(x,y)\,dy$.
\end{theorem}

% Source: Diego Torres Tejeda/Notes - 27.04 - Diego Torres Tejeda.pdf, p. 1
% Generator: Gemini 3.6 Flash (High)
\begin{proposition}[Cavalieri's principle]
\label{prop:cavalieri_principle}
Let $A, B \subseteq \mathbb{R}^{n+1}$ be Jordan measurable sets. For each height $t \in \mathbb{R}$, denote by $A_t := \{x \in \mathbb{R}^n \mid (x,t) \in A\}$ and $B_t := \{x \in \mathbb{R}^n \mid (x,t) \in B\}$ the $n$-dimensional cross-sections at $t$. If the cross-sectional volumes agree for all $t \in \mathbb{R}$, i.e.\
\[
\mu_n(A_t) = \mu_n(B_t) \quad \forall t \in \mathbb{R},
\]
then the total $(n+1)$-dimensional volumes are equal: $\mu_{n+1}(A) = \mu_{n+1}(B)$.
\end{proposition}

\begin{proof}
By Fubini's Theorem (\cref{thm:fubini}), integrating the indicator function over slices gives:
\[
\mu_{n+1}(A) = \int_{\mathbb{R}^{n+1}} \chi_A(x,t) \,dx\,dt = \int_{\mathbb{R}} \left( \int_{\mathbb{R}^n} \chi_{A_t}(x) \,dx \right) dt = \int_{\mathbb{R}} \mu_n(A_t) \,dt.
\]
Replacing $\mu_n(A_t)$ with $\mu_n(B_t)$ yields $\mu_{n+1}(B)$.
\end{proof}

\begin{theorem}[Volume of a cone over a measurable base]
\label{thm:cone_volume_formula}
Let $\Omega \subseteq \mathbb{R}^n$ be a bounded, Jordan measurable domain. The \newterm{cone over $\Omega$} \germanfor{Kegel \"uber $\Omega$} with vertex at $(0,1) \in \mathbb{R}^n \times \mathbb{R}$ and base $\Omega \times \{0\}$ is defined by
\[
C\Omega := \big\{(x, t) \in \mathbb{R}^n \times [0,1] \mid x \in (1-t)\Omega\big\}.
\]
Then $C\Omega$ is Jordan measurable in $\mathbb{R}^{n+1}$, and its volume is given by
\[
\mu_{n+1}(C\Omega) = \frac{\mu_n(\Omega)}{n+1}.
\]
\end{theorem}

\begin{proof}
For any $t \in [0,1]$, the cross-section of $C\Omega$ at height $t$ is $(C\Omega)_t = (1-t)\Omega \subseteq \mathbb{R}^n$. By the scaling property of $n$-dimensional Jordan measure, $\mu_n\big((1-t)\Omega\big) = (1-t)^n \mu_n(\Omega)$. Applying Cavalieri's Principle (\cref{prop:cavalieri_principle}):
\[
\mu_{n+1}(C\Omega) = \int_0^1 \mu_n\big((1-t)\Omega\big) \,dt = \mu_n(\Omega) \int_0^1 (1-t)^n \,dt = \mu_n(\Omega) \left[ -\frac{(1-t)^{n+1}}{n+1} \right]_0^1 = \frac{\mu_n(\Omega)}{n+1}.
\]
\end{proof}

\begin{center}
% Generator: Gemini 3.6 Flash (High) --- FIG-CONE: Cavalieri's principle for the volume of a cone over a domain
\begin{tikzpicture}[>=Latex, scale=1.0]

  % Axes
  \draw[->, thick] (-0.5,0) -- (3.2,0) node[right, font=\footnotesize] {$\mathbb{R}^n$};
  \draw[->, thick] (0,-0.3) -- (0,2.6) node[above, font=\footnotesize] {$t \in \mathbb{R}$};

  % Vertex at (0, 1.8)
  \fill[red!80!black] (0, 1.8) circle (1.8pt) node[left, font=\footnotesize] {$(0,1)$};

  % Base domain Omega at t=0
  \draw[very thick, orange!90!black] (0.8,0) -- (2.6,0);
  \fill[orange!90!black] (0.8,0) circle (1.5pt);
  \fill[orange!90!black] (2.6,0) circle (1.5pt);
  \node[orange!90!black, font=\footnotesize, below] at (1.7, -0.05) {$\Omega$};

  % Cone sides
  \draw[blue!75!black, thick] (0, 1.8) -- (0.8,0);
  \draw[blue!75!black, thick] (0, 1.8) -- (2.6,0);

  % Shaded region for C\Omega
  \fill[blue!10, opacity=0.7] (0, 1.8) -- (0.8,0) -- (2.6,0) -- cycle;
  \node[blue!80!black, font=\small] at (1.1, 0.7) {$C\Omega$};

  % Cross-section at height t
  \draw[dashed, gray!70!black] (-0.3, 1.0) -- (2.8, 1.0);
  \draw[very thick, green!50!black] (0.355, 1.0) -- (1.155, 1.0);
  \fill[green!50!black] (0.355, 1.0) circle (1.4pt);
  \fill[green!50!black] (1.155, 1.0) circle (1.4pt);
  \node[green!50!black, font=\scriptsize, above] at (0.75, 1.05) {$(1-t)\Omega$};

  \node[font=\scriptsize, left] at (0, 1.0) {$t$};
  \node[font=\scriptsize, left] at (0, 0.0) {$0$};

\end{tikzpicture}
\end{center}

% Source: Corsin Nick/Class Notes/Week 8.pdf, p. 8
\begin{example}[Fubini on the unit square]
\label{ex:fubini_unit_square}
\[
\begin{aligned}
\int_{[0,1]^2} x^{y+1}\,dx\,dy
  &= \int_0^1 dy \int_0^1 dx\ x^{y+1}
   = \int_0^1 dy\ \left[\frac{x^{y+2}}{y+2}\right]_0^1 \\
  &= \int_0^1 dy\ \frac{1}{y+2}
   = \log|y+2|\,\Big|_0^1
   = \log\!\left(\tfrac32\right).
\end{aligned}
\]
\end{example}

\begin{ainote}
Corsin writes the first line as $\int_0^1 dx \int_0^1 dy$, but the antiderivative taken on the
next line is in $x$, and it produces $x^{y+2}/(y+2)$, evaluated from $0$ to $1$. The
differentials are therefore swapped relative to the computation actually performed. The order is corrected
above and the result is unaffected.

The example is well chosen: integrating $x^{y+1}$ in $y$ first would mean integrating an
exponential in the exponent, which is far worse. Fubini's real use is that you may pick the
easier order, and here only one of the two is comfortable.
\end{ainote}

% Supplement: Sascha Brack/Class Notes/Week_08_Notes_Friday.pdf, pp. 4--5
\begin{example}[Changing the order of integration]
\label{ex:change_order_integration}
Consider the iterated integral
\[\int_0^2 \int_{y^2}^{2\sqrt{2y}} f(x,y) \,dx\,dy.\]
Let us rewrite it with the integration order reversed (first $y$, then $x$).

The region is defined by $0 \leq y \leq 2$ and $y^2 \leq x \leq 2\sqrt{2y}$.
Notice that for $y \in [0,2]$, the lower bound is $x \geq 0$ and the upper bound is $x \leq 2\sqrt{4} = 4$, so $x$ ranges from $0$ to $4$.
We now solve the inequalities for $y$:
$x \leq 2\sqrt{2y} \implies x^2 \leq 8y \implies y \geq \frac{1}{8}x^2$.
$x \geq y^2 \implies y \leq \sqrt{x}$ (since $x,y \geq 0$).
Thus, the new bounds are $0 \leq x \leq 4$ and $\frac{1}{8}x^2 \leq y \leq \sqrt{x}$. The reversed integral is:
\[\int_0^4 \int_{\frac{1}{8}x^2}^{\sqrt{x}} f(x,y) \,dy\,dx.\]
\exinfo{The same region, with the same two bounds, is Problem 14 of the Analysis II examination of
August 2025 (Prof.\ Laura Kobel-Keller), a box problem where only the reversed integral was
graded.}
\end{example}

% Generator: Claude Opus 5 (High) --- FIG-20-01, the two slicings of one region.
% Requested in feedback/brainstorming_improvements.md, which names swapping bounds over a
% non-rectangular domain as a major pain point and asks for exactly this picture. Drawn against
% the region of ex:change_order_integration rather than the brainstorm's triangle, because a
% triangle makes the two orders look symmetric and this region does not: the curved boundaries
% are what force the algebra in the example above.
%
% Arithmetic check. The two boundary curves are the SAME two curves, read two ways:
%   x = y^2        <=>  y = sqrt(x)      (upper boundary)
%   x = 2 sqrt(2y) <=>  y = x^2/8        (lower boundary)
% They meet where sqrt(x) = x^2/8, i.e. x^{3/2} = 8, x = 4, y = 2 -- and at the origin.
% Both panels must therefore show the SAME shaded region; only the slice differs.
%   LEFT (dx dy, outer y): horizontal slice at y = 1.2, running x = y^2 = 1.44
%       to x = 2 sqrt(2.4) = 3.0984.
%   RIGHT (dy dx, outer x): vertical slice at x = 2.2, running y = 2.2^2/8 = 0.605
%       to y = sqrt(2.2) = 1.4832.
% Both slice endpoints lie on the boundary curves by construction, which is the property the
% figure asserts and the one worth re-checking if the coordinates are ever touched.
\begin{center}
\begin{tikzpicture}[x=1.05cm, y=1.05cm, >=Latex,
    bdry/.style={blue!70!black, very thick},
    slice/.style={BrickRed, very thick, <->}]

  % ---------------- left panel: horizontal slices, the order as given ----------------
  \begin{scope}
    \fill[blue!8]
      plot[domain=0:4, samples=60] ({\x}, {sqrt(\x)}) --
      plot[domain=4:0, samples=60] ({\x}, {(\x)*(\x)/8}) -- cycle;

    \draw[->, gray!70] (-0.25,0) -- (4.9,0) node[right, font=\footnotesize, text=black] {$x$};
    \draw[->, gray!70] (0,-0.25) -- (0,2.7) node[above, font=\footnotesize, text=black] {$y$};

    \draw[bdry] plot[domain=0:4, samples=60] ({\x}, {sqrt(\x)});
    \draw[bdry] plot[domain=0:4, samples=60] ({\x}, {(\x)*(\x)/8});
    % Placed at x = 3.4 on the upper curve (y = sqrt(3.4) = 1.844) rather than near x = 1: the
    % dashed guide at y = 1.2 runs across the left half of this panel and the label collided
    % with it there.
    \node[above left, font=\scriptsize, blue!70!black, inner sep=1pt] at (3.4,1.844) {$x=y^2$};
    \node[below right, font=\scriptsize, blue!70!black, inner sep=1pt] at (3.05,1.16)
      {$x=2\sqrt{2y}$};

    \draw[slice] (1.44,1.2) -- (3.0984,1.2);
    \draw[gray!60, dashed] (0,1.2) -- (1.44,1.2);
    \node[left, font=\scriptsize] at (-0.1,1.2) {$y$};

    \fill (4,2) circle (1.4pt);
    \node[right, font=\scriptsize, inner sep=2pt] at (4.05,2.1) {$(4,2)$};

    \node[font=\scriptsize, BrickRed] (Lslab) at (1.05,2.5) {$y^2 \le x \le 2\sqrt{2y}$};
    \draw[BrickRed, thin] (Lslab) -- (1.9,1.28);
    \node[font=\footnotesize, align=center] at (2.0,-0.95)
      {$\displaystyle\int_0^2\!\!\int_{y^2}^{2\sqrt{2y}} f \,dx\,dy$};
    \node[font=\scriptsize, gray!50!black] at (2.0,-1.75) {fix $y$, sweep $x$};
  \end{scope}

  % ---------------- right panel: vertical slices, the order after swapping ------------
  \begin{scope}[shift={(6.6,0)}]
    \fill[blue!8]
      plot[domain=0:4, samples=60] ({\x}, {sqrt(\x)}) --
      plot[domain=4:0, samples=60] ({\x}, {(\x)*(\x)/8}) -- cycle;

    \draw[->, gray!70] (-0.25,0) -- (4.9,0) node[right, font=\footnotesize, text=black] {$x$};
    \draw[->, gray!70] (0,-0.25) -- (0,2.7) node[above, font=\footnotesize, text=black] {$y$};

    \draw[bdry] plot[domain=0:4, samples=60] ({\x}, {sqrt(\x)});
    \draw[bdry] plot[domain=0:4, samples=60] ({\x}, {(\x)*(\x)/8});
    \node[above left, font=\scriptsize, blue!70!black, inner sep=1pt] at (1.15,1.07)
      {$y=\sqrt{x}$};
    \node[below right, font=\scriptsize, blue!70!black, inner sep=1pt] at (3.05,1.16)
      {$y=x^2/8$};

    \draw[slice] (2.2,0.605) -- (2.2,1.4832);
    \draw[gray!60, dashed] (2.2,0) -- (2.2,0.605);
    \node[below, font=\scriptsize] at (2.2,-0.08) {$x$};

    \fill (4,2) circle (1.4pt);
    \node[right, font=\scriptsize, inner sep=2pt] at (4.05,2.1) {$(4,2)$};

    \node[font=\scriptsize, BrickRed] (Rslab) at (1.15,2.5) {$\tfrac{x^2}{8} \le y \le \sqrt{x}$};
    \draw[BrickRed, thin] (Rslab) -- (2.2,1.55);
    \node[font=\footnotesize, align=center] at (2.0,-0.95)
      {$\displaystyle\int_0^4\!\!\int_{x^2/8}^{\sqrt{x}} f \,dy\,dx$};
    \node[font=\scriptsize, gray!50!black] at (2.0,-1.75) {fix $x$, sweep $y$};
  \end{scope}
\end{tikzpicture}
\end{center}

\begin{remark}[Reading the swap off the picture]
\label{rem:fubini_swap_picture}
The region is the same in both panels; only the direction of the slice changes, and with it which
variable is the outer one. That is the whole of Fubini's theorem for a non-rectangular domain, and
the two mechanical rules follow from looking:

\begin{itemize}
  \item \textbf{The outer limits are constants, and they are the shadow of the region on the outer
    axis.} On the left the outer variable is $y$ and the region's shadow on the $y$-axis is
    $[0,2]$; on the right the outer variable is $x$ and its shadow is $[0,4]$. The two are
    different numbers, which is the most common slip.
  \item \textbf{The inner limits are functions of the outer variable, read off where the slice
    enters and leaves the region.} Entering and leaving happens on the \emph{same two curves} in
    both panels --- there is only one pair of boundary curves here. What changes is which of
    $x = y^2$ and $y = \sqrt{x}$ is the useful way to write the first of them.
\end{itemize}

The reason the algebra in the example above is not merely a relabelling is the second bullet: each
boundary has to be re-solved for the new inner variable. A region whose boundary cannot be solved
for one of the variables cannot be swapped at all without first splitting it, and that is the
situation \cref{ex:change_of_variables_domain} runs into.
\end{remark}

\begin{example}[Surface parametrization and Fubini]
\label{ex:surface_parametrization_integral}
Let $f(u,v) := (u\cos v, u\sin v, u^2)\transp \in \mathbb{R}^3$ for $(u,v) \in \mathbb{R}^2$. Write the third coordinate $z(u,v)$ as a function of the first two, $\tilde{z}(x,y)$, and compute $\int_0^y \int_0^x \tilde{z}(x',y') \,dx'\,dy'$.

We have $x(u,v) = u\cos v$ and $y(u,v) = u\sin v$. Then $x^2 + y^2 = u^2\cos^2 v + u^2\sin^2 v = u^2 = z(u,v)$.
So, we define $\tilde{z}(x,y) := x^2 + y^2$.
The integral is:
\begin{align}
\int_0^y \int_0^x (x'^2 + y'^2) \,dx'\,dy'
&= \int_0^y \left[ \frac{x'^3}{3} + x'y'^2 \right]_{x'=0}^{x'=x} \,dy' \nonumber \\
&= \int_0^y \left( \frac{x^3}{3} + x y'^2 \right) \,dy' \nonumber \\
&= \left[ \frac{x^3 y'}{3} + x \frac{y'^3}{3} \right]_{y'=0}^{y'=y} \nonumber \\
&= \frac{xy}{3}(x^2 + y^2). \label{eq:surface_integral_result}
\end{align}
\end{example}

% Source: Corsin Nick/Class Notes/Week 8.pdf, pp. 8--10
\begin{example}[A domain that Fubini alone cannot handle]
\label{ex:change_of_variables_domain}
Consider
\[\int_A \frac{y}{\sqrt{x}}\,dy\,dx, \qquad
A := \left\{(x,y) \in \mathbb{R}^2 \ :\ x,y>0,\ \ 1 \leq \frac{y}{\sqrt x} \leq 2,\ \ 1 \leq xy \leq 2\right\}.\]
This is not easily integrated with Fubini alone. But we can apply a transformation of variables
to simplify the integration domain $A$: take
\[\Phi : U \to V, \qquad \begin{pmatrix}x\\y\end{pmatrix} \mapsto
  \begin{pmatrix}u(x,y)\\v(x,y)\end{pmatrix} = \begin{pmatrix}\tfrac{y}{\sqrt x}\\[2pt] xy\end{pmatrix},\]
with $U$ and $V$ still to be determined, subject only to their being open.

\textbf{Check injectivity.} Suppose, for $x,y,x',y' > 0$,
\[\text{\textbf{(1)}}\ \ \frac{y}{\sqrt x} = \frac{y'}{\sqrt{x'}},
\qquad\qquad \text{\textbf{(2)}}\ \ xy = x'y'.\]
From \textbf{(2)}, $\dfrac{x}{x'} = \dfrac{y'}{y}$. Inserting into \textbf{(1)},
\[\frac{1}{\sqrt x} = \frac{y'}{y}\cdot\frac{1}{\sqrt{x'}} = \frac{x}{x'}\cdot\frac{1}{\sqrt{x'}}
\quad\Longrightarrow\quad x^{3/2} = (x')^{3/2} \quad\Longrightarrow\quad x = x',\]
and consequently $\dfrac{y'}{y} = \dfrac{x}{x'} = 1$. So, for $U := (0,\infty)^2$ and
$V := \Phi(U)$, the map $\Phi$ is a bijection between open sets.

\textbf{Functional determinant.} We calculate $|\det \Jac\Phi|$; if it is strictly positive we
may proceed, since $\Phi$ is then a local diffeomorphism as well as a bijection, hence a
diffeomorphism. Now
\[\Jac\Phi(x,y) = \begin{pmatrix} -\dfrac{y}{2\sqrt{x}^{\,3}} & \dfrac{1}{\sqrt x} \\[2ex] y & x \end{pmatrix},\]
so that
\[\big|\det \Jac\Phi(x,y)\big| = \left|-\frac32\,\frac{y}{\sqrt x}\right| = \frac32\,\frac{y}{\sqrt x} = \frac32\,u(x,y).\]
So, $\Phi$ is a diffeomorphism and $du\,dv = \tfrac32 u\,dx\,dy$. Furthermore, by design,
\[\Phi(A) = (1,2)^2.\]
So, with the change of variables (\cref{thm:change_of_variables}), followed by Fubini
(\cref{thm:fubini}),
\[\int_A \frac{y}{\sqrt x}\,dy\,dx
  = \int_{(1,2)^2} u\,\frac{1}{\tfrac32 u}\,du\,dv
  = \frac23\int_1^2 dv \int_1^2 du
  = \frac23.\]
\end{example}

\begin{ainote}
The design principle is worth stating, because it is the reverse of how substitution is usually
taught. Here $\Phi$ was not chosen to simplify the \emph{integrand} but to straighten the
\emph{domain}: the two constraints defining $A$ are literally $1 \leq u \leq 2$ and
$1 \leq v \leq 2$, so $\Phi$ turns an awkward curvilinear region into a square, and only then
does Fubini become available. That the integrand $y/\sqrt x$ is itself $u$, and cancels against
most of the Jacobian factor, is a bonus of the same choice.
\end{ainote}

\begin{aiexercise}[An order of integration that cannot be done]
% Generator: Claude Opus 5 (High)
\label{ex:ai_order_swap_gaussian}
Consider
\[I := \int_0^1\!\!\int_x^1 e^{y^2}\,dy\,dx.\]
\begin{enumerate}[label=\textbf{(\alph*)}]
  \item Explain what goes wrong if one attempts the inner integral as written.
  \item Describe the region of integration by inequalities, and rewrite $I$ with the two
    integrations in the opposite order.
  \item Evaluate $I$.
\end{enumerate}
\end{aiexercise}

% Source: exercises_2024/mock.pdf (Analysis II mock examination, Spring Semester 2024,
% Prof. Joaquim Serra), Exercise 12, p. 3
% Extractor: Claude Opus 5 (High)
% Filed with Fubini rather than with the improper integrals of chapter 19 because the whole
% problem is the product splitting; the one-dimensional pieces are then immediate. Part (a) is the
% trap and is what earns the block its place: the printed value sqrt(pi) is the answer to a
% one-dimensional Gaussian, and alpha = 0 is exactly the case where the z-integral does not exist.
\begin{exercise}[True or False \textnormal{(a Gaussian slab, and what happens at $\alpha=0$)}]
\label{ex:gaussian_slab_alpha_asymptotics}
For $\alpha \geq 0$ consider
\[I_\alpha := \int_{\mathbb{R}^3} e^{-x^2-y^2-\alpha\lvert z\rvert}\,dx\,dy\,dz .\]
Decide whether each of the following is true or false, with justification.
\begin{enumerate}[label=\textbf{(\alph*)}]
  \item $I_0 = \sqrt{\pi}$.
  \item $\alpha I_\alpha = I_1$.
  \item $\displaystyle\lim_{\alpha\to+\infty}\big(\alpha+\sqrt{\alpha}\big)I_\alpha = 2\pi$.
\end{enumerate}
\exinfo{This exercise is taken from the mock examination for Analysis II, Spring Semester 2024
(Prof.\ Joaquim Serra), where it appears as the twelfth block of multiple-choice questions. Its
paper supplies $\int_{\mathbb{R}}e^{-t^2}\,dt = \sqrt{\pi}$ as a formula that may be used without
proof; \cref{ex:gaussian_integral_polar} derives it.}
\end{exercise}

% Sonnet 5 (Medium)
We now want to look at a powerful new integration technique using ODEs.

% Originally: content/week-08.tex, \section{Feynman's trick} (Corsin Week 8)

% Source: Corsin Nick/Class Notes/Week 8.pdf, p. 11
\section{Feynman's trick}
\label{sec:feynmans_trick}

\begin{theorem}[Differentiation under the integral sign]
\label{thm:feynmans_trick}
Let $U \subseteq \mathbb{R}^n\times\mathbb{R}$ be open, $f \in C^1(U)$, and let
$K \subseteq \mathbb{R}^n$ be compact such that $K\times[a,b] \subseteq U$. Then the function
\[(a,b) \to \mathbb{R}, \qquad y \mapsto \int_K f(x,y)\,dx\]
is continuously differentiable, and
\[\frac{d}{dy}\int_K f(x,y)\,dx = \int_K \frac{\partial}{\partial y}f(x,y)\,dx
\qquad \text{for all } y \in (a,b).\]
\end{theorem}

% Source: Corsin Nick/Class Notes/Week 8.pdf, pp. 11--13
\begin{example}[$\int_0^1 \frac{\log(1+x)}{1+x^2}\,dx$ by Feynman's trick]
\label{ex:feynman_log_integral}
We want to compute the integral
\[I(1) = \int_0^1 \frac{\log(1+x)}{1+x^2}\,dx,\]
which has no convenient antiderivative. By introducing a parameter $\alpha \in (-1,1)$:
\[I(\alpha) := \int_0^1 \frac{\log(1+\alpha x)}{1+x^2}\,dx.\]
Note that the integrand is a $C^1$ function on
$U := \big(-\tfrac{1}{1+\varepsilon},\ \tfrac1\varepsilon\big) \times (-\varepsilon,\ 1+\varepsilon)$
in the variables $(x,\alpha)$, with $K := [0,1]$ and $[a,b] := [0,1]$ satisfying
$K\times[a,b] \subseteq U$ for $0 < \varepsilon < 1$ (pay attention to the logarithm, which is
what forces $1+\alpha x > 0$). So, we may differentiate under the integral!

\textbf{We obtain an ordinary differential equation.} By \cref{thm:feynmans_trick},
\[I'(\alpha) = \frac{d}{d\alpha}\int_0^1 \frac{\log(1+\alpha x)}{1+x^2}\,dx
  = \int_0^1 \frac{x}{(1+\alpha x)(1+x^2)}\,dx.\]
Partial fractions: writing
\[\frac{x}{(1+\alpha x)(1+x^2)} = \frac{A}{1+\alpha x} + \frac{Bx+C}{1+x^2}\]
and clearing denominators gives $A + Ax^2 + Bx + B\alpha x^2 + C + C\alpha x = x$, so comparing
coefficients,
\[-A = C, \qquad -\tfrac{A}{\alpha} = B, \qquad B + C\alpha = 1
  \implies -\tfrac{A}{\alpha} - A\alpha = 1,\]
whence
\[A = -\frac{\alpha}{1+\alpha^2}, \qquad B = \frac{1}{1+\alpha^2}, \qquad C = \frac{\alpha}{1+\alpha^2}.\]
Therefore
\[
\begin{aligned}
I'(\alpha) &= \left(\int_0^1 \frac{-\alpha}{1+\alpha x} + \frac{x+\alpha}{1+x^2}\,dx\right)\frac{1}{1+\alpha^2} \\
&= \frac{1}{1+\alpha^2}\Big[-\log(1+\alpha x) + \tfrac12\log(1+x^2) + \alpha\arctan(x)\Big]_0^1 \\
&= \frac{1}{1+\alpha^2}\Big[\log\sqrt2 + \tfrac{\alpha\pi}{4} - \log(1+\alpha)\Big].
\end{aligned}
\]

\textbf{With initial value} $I(0) = \int_0^1 \frac{\log(1)}{1+x^2}\,dx = 0$, we obtain from the
fundamental theorem of calculus
\[
\begin{aligned}
\int_0^1 \frac{\log(1+x)}{1+x^2}\,dx = I(1) &= I(1)-I(0) \\
&= \int_0^1\left(\frac{\log\sqrt2}{1+\alpha^2} + \frac{\pi}{4}\,\frac{\alpha}{1+\alpha^2}
   - \underbrace{\frac{\log(1+\alpha)}{1+\alpha^2}}_{\textstyle \to\, I(1)}\right)d\alpha.
\end{aligned}
\]
The last term is $I(1)$ again, the integral we started from with $\alpha$ merely renamed, so
bringing it to the left-hand side,
\[
\begin{aligned}
2\,I(1) &= \Big[\log\sqrt2\,\arctan(\alpha) + \tfrac{\pi}{8}\log(1+\alpha^2)\Big]_0^1
 = \frac{\pi\log\sqrt2}{4} + \frac{\pi\log\sqrt2}{4}, \\[1ex]
\implies\quad I(1) &= \frac{\pi\log\sqrt2}{4} = \boxed{\ \frac{\pi\log 2}{8}\ }.
\end{aligned}
\]
\end{example}

% Source: Jérôme Paschoud/Notizen/Woche 9.1 Variablenwechsel.pdf, p. 4
% Generator: Gemini 3.1 Pro (High)
\begin{example}[Dirichlet integral via Feynman's trick]
\label{ex:feynman_dirichlet_integral}
Consider the integral $\int_0^\infty \frac{\sin x}{x} \,dx$. We can evaluate it by introducing a damping factor $e^{-tx}$ and defining
\[I(t) := \int_0^\infty e^{-tx} \frac{\sin x}{x} \,dx \qquad \text{for } t > 0.\]
Differentiating under the integral sign with respect to $t$ (which is justified since the integrand and its derivative decay rapidly for $t>0$):
\[I'(t) = \int_0^\infty \frac{\partial}{\partial t} \left( e^{-tx} \frac{\sin x}{x} \right) \,dx = \int_0^\infty -x e^{-tx} \frac{\sin x}{x} \,dx = -\int_0^\infty e^{-tx} \sin x \,dx.\]
This standard integral can be evaluated using integration by parts twice, or by using Euler's formula $e^{ix} = \cos x + i\sin x$ and noting that $e^{-tx}\sin x = \operatorname{Im}(e^{(-t+i)x})$:
\[\int_0^\infty e^{(-t+i)x} \,dx = \left[ \frac{e^{(-t+i)x}}{-t+i} \right]_0^\infty = 0 - \frac{1}{-t+i} = \frac{t+i}{t^2+1}.\]
Taking the imaginary part gives $\int_0^\infty e^{-tx} \sin x \,dx = \frac{1}{t^2+1}$. Thus,
\[I'(t) = -\frac{1}{t^2+1} \implies I(t) = -\arctan(t) + C.\]
Since $\lim_{t \to \infty} I(t) = 0$, we must have $C = \lim_{t \to \infty} \arctan(t) = \frac{\pi}{2}$.
Therefore, $I(t) = \frac{\pi}{2} - \arctan(t)$.
The original integral corresponds to the limit as $t \to 0^+$ (which requires an extra justification step usually provided by Abel's theorem, but we formally take the limit):
\[\int_0^\infty \frac{\sin x}{x} \,dx = \lim_{t \to 0^+} \left( \frac{\pi}{2} - \arctan(t) \right) = \frac{\pi}{2}.\]
\end{example}

\begin{ainote}
Corsin's final integration writes the $\alpha$-term as
$\tfrac{\pi}{2}\cdot\tfrac{\alpha}{1+\alpha^2}$, integrating to $\tfrac{\pi}{4}\log(1+\alpha^2)$.
Both are a factor $2$ away from what $I'(\alpha)$ gives: his own line above has
$\tfrac{\alpha\pi}{4}$, so the term is $\tfrac{\pi}{4}\cdot\tfrac{\alpha}{1+\alpha^2}$ and its
antiderivative is $\tfrac{\pi}{8}\log(1+\alpha^2)$. Corrected above, and the correction is
self-checking, since it is exactly what makes his two bracket terms come out equal, as his own
next line asserts:
$\tfrac{\pi}{8}\log 2 = \tfrac{\pi}{8}\cdot 2\log\sqrt2 = \tfrac{\pi\log\sqrt2}{4}$. His final
answer $\tfrac{\pi\log 2}{8}$ is correct and matches the known value.
\end{ainote}

% Generator: Claude Opus 5 (Medium)
\begin{ainote}[Why this is called a trick, and when to reach for it]
The move that makes this work is not the differentiation. It is noticing that the
\emph{antiderivative} of the resulting integrand is elementary while the original one is not.
The pattern is worth naming:
\begin{enumerate}[label=\textbf{(\alph*)}]
  \item Insert a parameter $\alpha$ so that the target is $I(\alpha_1)$ for some
    $\alpha_1$, and so that $I(\alpha_0)$ is trivial for some other $\alpha_0$. Here
    $I(0)=0$, because $\log 1 = 0$.
  \item Differentiate under the integral. The $\log$ disappears, leaving a rational function
    that partial fractions can handle.
  \item Integrate $I'$ back from $\alpha_0$ to $\alpha_1$.
\end{enumerate}
The self-reference in step 3, where the answer reappears on the right-hand side so that one
solves for $I(1)$ rather than computing it, is specific to this integral and is the reason it
is a classic. Do not expect it in general. Usually step 3 is an honest integration.
\end{ainote}

\begin{aiexercise}[Frullani's integral]
% Generator: Claude Opus 5 (High)
\label{ex:ai_frullani}
Fix $0 < a < b$. The aim is to show that
\[\int_0^\infty \frac{e^{-ax} - e^{-bx}}{x}\,dx = \log\!\left(\frac{b}{a}\right),\]
an answer that depends on $a$ and $b$ only through their ratio.
\begin{enumerate}[label=\textbf{(\alph*)}]
  \item The integrand is undefined at $x = 0$. Compute $\lim_{x\to0^+}$ of it, so that the
    integrand extends continuously to $[0,\infty)$.
  \item Hold $a$ fixed and let the second exponent be the parameter: define
    \[I(t) := \int_0^\infty \frac{e^{-ax} - e^{-tx}}{x}\,dx \qquad \text{for } t > 0.\]
    Differentiate under the integral sign and evaluate the resulting integral, which is
    elementary.
  \item Identify the value of $t$ for which $I(t)$ is trivially zero, and integrate $I'$ from
    there to $t = b$.
\end{enumerate}
\end{aiexercise}

% Generator: Claude Opus 5 (High)
% Section added 2026-08-11. Both exercises are transcribed from the Probeprfg mock exams (see the
% % Source: comments below); the definition, the connecting prose and the remarks are written for
% this document, which had no notion of a centroid anywhere. The two problems come from different
% papers and are placed together because each is the other's payoff: one computes the centroid,
% the other says what the centroid is for.

\section{Centroids}
\label{sec:centroids}

% Transition: Claude Opus 5 (High)
Every integral in this chapter so far has been computed for its own sake. The simplest thing one
actually does with an integral over a region is average a coordinate over it, and the resulting
point is worth a name.

\begin{definition}[Centroid of a plane region]
\label{def:centroid}
Let $B \subseteq \mathbb{R}^2$ be Jordan measurable with $\vol_2(B) > 0$. The \newterm{centroid}
\germanfor{Fl\"achenschwerpunkt} of $B$ is the point $z_B := (x_B, y_B)$ with
\[x_B := \frac{1}{\vol_2(B)}\int_B x \,d\vol_2(x,y), \qquad
  y_B := \frac{1}{\vol_2(B)}\int_B y \,d\vol_2(x,y).\]
\end{definition}

The same formula defines the centroid of a Jordan measurable set in any $\mathbb{R}^n$, one
coordinate at a time; nothing below uses the restriction to the plane except the triangle.

\begin{remark}[Why the normalisation is there]
Without the factor $1/\vol_2(B)$ the two integrals would scale with the size of $B$ as well as
with its position, and doubling a region would move its \qt{centre}. Dividing by the area makes
$z_B$ an \emph{average} of the coordinate functions rather than a total, and the immediate
consequence is that $z_B$ is where one expects it to be: for a region symmetric about a line, the
centroid lies on that line, because the coordinate perpendicular to it integrates to zero.
\end{remark}

\begin{center}
% Generator: Claude Opus 5 (High) --- FIG-20-02: the two halves of the exercise below, drawn.
% This section defines a *point* and then asserts two things about where it sits, with no picture
% anywhere: the reader meets the weighted-average formula of part (a) as an identity between two
% integrals and never sees that it puts $z_B$ on the segment joining the two sub-centroids.
% Both panels are exact, not schematic --- the coordinates below are the computed centroids, so
% the picture can be measured against the formulas it illustrates.
% `lbl' backs a label with a soft white plate so it stays readable where it lands on a shaded
% region or a box edge -- the same idiom as the surface figure in ch. 21, and the thing that
% decides whether these read on a grey e-reader screen rather than only at 300dpi in colour.
\begin{tikzpicture}[>=Latex, scale=1.0,
  lbl/.style={fill=white, fill opacity=0.8, text opacity=1, inner sep=1.2pt,
              rounded corners=1pt}]

  % ---------- Left: centroids add, as an area-weighted average ----------
  % B_1 = [0,2]x[0,1.8], area 3.6, centroid (1, 0.9)
  % B_2 = [2,3]x[0,0.9], area 0.9, centroid (2.5, 0.45)
  % weight of B_2 is 0.9/4.5 = 1/5 exactly, so
  % z_B = z_1 + (1/5)(z_2 - z_1) = (1.3, 0.81)
  \node[font=\small\bfseries, text=TextBoldColor] at (1.5, 3.0) {Centroids add};

  \fill[blue!10, opacity=0.7] (0,0) rectangle (2,1.8);
  \fill[orange!12, opacity=0.7] (2,0) rectangle (3,0.9);
  \draw[blue!75!black, thick] (0,0) rectangle (2,1.8);
  \draw[orange!90!black, thick] (2,0) rectangle (3,0.9);
  \node[blue!80!black, font=\scriptsize] at (0.45, 1.5) {$B_1$};
  \node[orange!90!black, font=\scriptsize] at (2.75, 0.68) {$B_2$};

  % the segment joining the two sub-centroids, and z_B on it
  \draw[gray!70!black, dashed] (1,0.9) -- (2.5,0.45);
  \fill[blue!75!black] (1,0.9) circle (1.5pt)
    node[lbl, above left=-2pt, font=\scriptsize] {$z_{B_1}$};
  \fill[orange!90!black] (2.5,0.45) circle (1.5pt)
    node[lbl, below right=-2pt, font=\scriptsize] {$z_{B_2}$};
  \fill[red!80!black] (1.3,0.81) circle (1.8pt);
  % text=, not a bare colour, which would override the fill=white in `lbl'.
  \node[lbl, text=red!80!black, font=\scriptsize, above right=-2pt] at (1.3,0.81) {$z_B$};

  \node[font=\scriptsize, align=center] at (1.5,-0.75)
    {$\vol_2(B_2)/\vol_2(B) = \tfrac15$, and $z_B$ sits\\one fifth of the way along};

  % ---------- Right: the centroid of a triangle ----------
  % o=(0,0), p=(3,0), q=(1.2,2.4); centroid = (4.2/3, 2.4/3) = (1.4, 0.8)
  \begin{scope}[xshift=5.0cm]
    \node[font=\small\bfseries, text=TextBoldColor] at (1.5, 3.0) {\dots\ and a triangle's};

    \fill[blue!10, opacity=0.7] (0,0) -- (3,0) -- (1.2,2.4) -- cycle;
    \draw[blue!75!black, thick] (0,0) -- (3,0) -- (1.2,2.4) -- cycle;

    % the three medians, which meet at the centroid
    \draw[gray!70!black, thin] (0,0) -- (2.1,1.2);
    \draw[gray!70!black, thin] (3,0) -- (0.6,1.2);
    \draw[gray!70!black, thin] (1.2,2.4) -- (1.5,0);

    % midpoints
    \fill[gray!70!black] (2.1,1.2) circle (1.0pt);
    \fill[gray!70!black] (0.6,1.2) circle (1.0pt);
    \fill[gray!70!black] (1.5,0) circle (1.0pt);

    \fill[black] (0,0) circle (1.5pt) node[below left=-2pt, font=\scriptsize] {$o$};
    \fill[black] (3,0) circle (1.5pt) node[below right=-2pt, font=\scriptsize] {$p$};
    \fill[black] (1.2,2.4) circle (1.5pt) node[above=1pt, font=\scriptsize] {$q$};

    \fill[red!80!black] (1.4,0.8) circle (1.8pt);
    \node[lbl, text=red!80!black, font=\scriptsize, right=1pt] at (1.4,0.8) {$z_B$};

    \node[font=\scriptsize, align=center] at (1.5,-0.75)
      {the three medians meet at\\$z_B = \tfrac13(o+p+q)$};
  \end{scope}

\end{tikzpicture}
\end{center}

% Source: old_exams/fs2023/Probeprfg3.pdf (Probeprüfung Analysis II, 31 January 2023),
% Aufgabe 11, p. 4
% Extractor: Claude Opus 5 (High)
% The paper names no lecturer; see old-exam-mining.md.
\begin{exercise}[Centroids add, and the centroid of a triangle]
\label{ex:centroid_additivity_and_triangle}
Let $B \subseteq \mathbb{R}^2$ be Jordan measurable with $\vol_2(B) > 0$.
\begin{enumerate}[label=\textbf{(\alph*)}]
  \item Suppose $B = \bigcup_{i=1}^n B_i$ for pairwise disjoint Jordan measurable sets
    $B_1,\dots,B_n \subseteq \mathbb{R}^2$, each with $\vol_2(B_i) > 0$. Show that
    \[z_B = \sum_{i=1}^n \frac{\vol_2(B_i)}{\vol_2(B)}\,z_{B_i}.\]
  \item Show that if $B$ is a triangle with vertices $o = (0,0)$, $p = (r,0)$ and $q = (s,t)$,
    where $r,s,t > 0$, then $z_B = \tfrac13(o+p+q)$.
\end{enumerate}
\exinfo{This exercise is taken from the mock examination \emph{Probepr\"ufung Analysis II} of
31 January 2023, Problem 11. That paper names no lecturer.}
\exsol{sol:centroid_additivity_and_triangle}
\end{exercise}

% Source: old_exams/fs2023/Probeprfg2.pdf (Probeprüfung Analysis II, 17 August 2022, revised),
% Aufgabe 1, p. 1
% Extractor: Claude Opus 5 (High)
% The paper names no lecturer; see old-exam-mining.md. The original asks only for the minimising
% pair; the two-part split, and part (b) in particular, are added here, because the general
% statement is what makes the computation worth doing and the exam's own solution differentiates
% under the integral sign without ever saying that the answer is the centroid.
\begin{exercise}[What the centroid minimises]
\label{ex:centroid_minimises_second_moment}
Let $B := \{(x,y) \in \mathbb{R}^2 \mid x \geq 0,\ y \geq 0,\ x^2+y^2 \leq 1\}$ be the closed
quarter disc, and for $(s,t) \in \mathbb{R}^2$ put
\[I(s,t) := \int_B (x-s)^2 + (y-t)^2 \,d\vol_2(x,y).\]
\begin{enumerate}[label=\textbf{(\alph*)}]
  \item Determine the pair $(s,t)$ at which $I$ is minimal.
  \item Show that the answer to part \textbf{(a)} is no accident: for \emph{any} Jordan
    measurable $B \subseteq \mathbb{R}^2$ with $\vol_2(B) > 0$, the function $I$ is minimised
    exactly at $(s,t) = z_B$, and at no other point.
\end{enumerate}
\exinfo{Part \textbf{(a)} is taken from the mock examination \emph{Probepr\"ufung Analysis II} of
17 August 2022 (revised), Problem 1; that paper names no lecturer. Part \textbf{(b)} is added
here, since the general statement is what the computation is really about.}
\exsol{sol:centroid_minimises_second_moment}
\end{exercise}

% Originally: content/week-08.tex had a \section{Solutions} holding nothing but a placeholder
% note, which had already gone stale, since the statements were in fact transcribed. So 8.5
% stood in the document unsolved. Written here for the first time.
\section{Solutions}
\label{sec:change_of_variables_solutions}

\begin{exercisesolution}[Solution to \cref{ex:8.5}]
% Generator: Claude Opus 5 (High)
In each part the work is the same two steps: find $\lvert\det\Jac\rvert$, then translate every
constraint. As the exercise says, no integral is actually evaluated.
\begin{enumerate}[label=\textbf{\arabic*.}]
  \item \textbf{Polar coordinates.} Here $\Phi(r,\theta) := (r\cos\theta, r\sin\theta)$ has
    $\det\Jac\Phi = r > 0$ (\cref{ex:polar_spherical_coordinates}), so
    $dx_1dx_2 = r\,dr\,d\theta$. The three constraints translate one at a time:
    \begin{itemize}
      \item $1 < x_1^2+x_2^2 < 4$ becomes $1 < r < 2$;
      \item $x_1 > 0$ becomes $\cos\theta > 0$, i.e.\ $\theta \in (-\tfrac\pi2, \tfrac\pi2)$;
      \item $x_2 < x_1$ becomes $r\sin\theta < r\cos\theta$; since $r > 0$ and $\cos\theta > 0$
        this is $\tan\theta < 1$, i.e.\ $\theta < \tfrac\pi4$.
    \end{itemize}
    So the transformed domain is the rectangle
    $(r,\theta) \in (1,2)\times\big(-\tfrac\pi2,\tfrac\pi4\big)$ and
    \[\int_A x_1^2\sin(x_2)\,dx_1dx_2
      = \int_{1}^{2}\!\!\int_{-\pi/2}^{\pi/4} r^3\cos^2\theta\,\sin(r\sin\theta)\,d\theta\,dr,\]
    the extra power of $r$ coming from $x_1^2 = r^2\cos^2\theta$ times the Jacobian factor $r$.
  \item \textbf{$u := xy$, $v := x^2$.} Follow Sascha Brack's hint and fix the signs first. On $B$
    the constraint $x^2 < y$ forces $y > 0$, and then $xy > 1 > 0$ forces $x > 0$; so
    $u > 0$ and $v > 0$ throughout, and $x = \sqrt v$, $y = u/\sqrt v$.

    Differentiating in the convenient direction,
    \[\frac{\partial(u,v)}{\partial(x,y)}
      = \det\begin{pmatrix} y & x \\ 2x & 0\end{pmatrix} = -2x^2 = -2v,
      \qquad\text{so}\qquad dx\,dy = \frac{du\,dv}{2v}.\]
    Now the constraints. $1 < xy < 2$ becomes $1 < u < 2$ at once. The second is the interesting
    one: $x^2 < y < 2x^2$ reads $v < u/\sqrt v < 2v$, i.e.
    \[v^{3/2} < u < 2\,v^{3/2}, \qquad\text{equivalently}\qquad
      \Big(\frac{u}{2}\Big)^{2/3} < v < u^{2/3}.\]
    The integrand becomes $y^2e^{-xy} = (u^2/v)\,e^{-u}$, so
    \[\int_B y^2e^{-xy}\,dx\,dy
      = \int_{1}^{2}\!\!\int_{(u/2)^{2/3}}^{u^{2/3}} \frac{u^2e^{-u}}{2v^2}\,dv\,du.\]
    Note what did \emph{not} happen: straightening $1 < xy < 2$ into a pair of constant bounds
    did nothing for the other constraint, which stays curved. Straightening one constraint does
    not oblige the rest to cooperate.
  \item \textbf{A linear change of variables.} With $u := z$, $v := z-2y$, $w := 3x+y+z$,
    \[\frac{\partial(u,v,w)}{\partial(x,y,z)}
      = \det\begin{pmatrix} 0 & 0 & 1 \\ 0 & -2 & 1 \\ 3 & 1 & 1\end{pmatrix} = 6,
      \qquad\text{so}\qquad dx\,dy\,dz = \frac{du\,dv\,dw}{6},\]
    a constant, since the map is linear. Each constraint already \emph{is} one of the new
    coordinates, so the transformed domain is literally the box
    \[(u,v,w) \in (0,1)\times(1,2)\times(-2,0).\]
    Inverting the substitution gives $z = u$, $y = \tfrac{u-v}{2}$ and
    $x = \tfrac{1}{6}\big(2w - 3u + v\big)$, so $xyz = \tfrac{1}{12}\,u\,(u-v)\,(2w-3u+v)$ and
    \[\int_C xyz\,dx\,dy\,dz
      = \int_{0}^{1}\!\!\int_{1}^{2}\!\!\int_{-2}^{0}
        \frac{u\,(u-v)\,(2w-3u+v)}{72}\,dw\,dv\,du.\]
\end{enumerate}

\begin{ainote}
The three parts are graded deliberately, and comparing them is the lesson. In part 3 the change
of variables was \emph{chosen} to be the constraints, so the domain is a box and the Jacobian a
constant, and nothing is left to do. In part 1 the domain is a box in the new coordinates but the
Jacobian is not constant. In part 2 neither holds: the Jacobian varies and one of the two
constraints stays curved. Reading the substitution off the constraints, as in part 3, is the move
worth remembering, and it is available far more often than it looks.
\end{ainote}
\end{exercisesolution}

\begin{exercisesolution}[Solution to \cref{ex:ai_triangle_substitution}]
% Generator: Claude Opus 5 (High)
\textbf{1. The new domain and the Jacobian factor.} The map $\Psi(x,y) := (y-x,\ y+x)$ is
linear and invertible, with
\[\Psi^{-1}(u,v) = (x,y) = \left(\frac{v-u}{2},\ \frac{u+v}{2}\right),
  \qquad
  \det\Jac\Psi^{-1} = \det\begin{pmatrix} -1/2 & 1/2 \\ 1/2 & 1/2\end{pmatrix} = -\frac12,\]
so $dx\,dy = \tfrac12\,du\,dv$. The three sides of $T$ translate one at a time: $x \geq 0$ reads
$u \leq v$, the condition $y \geq 0$ reads $u \geq -v$, and $x+y \leq 1$ reads $v \leq 1$. Taking
strict inequalities for the interior gives exactly $0 < v < 1$ and $-v < u < v$, a triangle
again, but one whose horizontal slices run in the direction the integrand likes.

\textbf{2. The value.} On that region the integrand is $e^{u/v}$, and $v$ is constant along the
inner integration, so the inner integral is an exponential in $u$:
\[I = \frac12\int_0^1\!\!\int_{-v}^{v} e^{u/v}\,du\,dv
    = \frac12\int_0^1 \Big[v\,e^{u/v}\Big]_{u=-v}^{u=v}\,dv
    = \frac{e - e^{-1}}{2}\int_0^1 v\,dv
    = \frac{e - e^{-1}}{4}.\]

\textbf{3. The origin.} A single point is a Jordan null set, so removing it changes neither the
Jordan measurability of $T$ nor the value of the integral. What does need checking is that the
integral exists at all, and it does, because the integrand is \emph{bounded}: on $T$ one has
$\lvert y-x\rvert \leq y+x$, so the exponent $u/v$ lies in $[-1,1]$ and the integrand lies
between $e^{-1}$ and $e$. It is the boundedness, and not the smallness of the bad set on its
own, that rescues the computation.

\begin{remark}[Two ways to choose a substitution]
\label{rem:substitution_from_the_integrand}
\Cref{ex:change_of_variables_domain} chose $\Phi$ to straighten the \emph{domain}. Here the
choice comes from the \emph{integrand}, which is a function of $y-x$ and $y+x$ and of nothing
else. That the domain also becomes pleasant is a bonus, and a mild one: the new region is still
a triangle. What made the integral collapse is that the inner integration now runs along the
direction in which the integrand is a plain exponential.
\end{remark}
\end{exercisesolution}

\begin{exercisesolution}[Solution to \cref{ex:ai_order_swap_gaussian}]
% Generator: Claude Opus 5 (High)
\textbf{1.} The inner integral asks for an antiderivative of $y \mapsto e^{y^2}$, and there is
none in elementary terms. Nothing about the outer integral repairs this, so the order as written
is a dead end.

\textbf{2.} The region is the triangle below the diagonal, and it admits two descriptions:
\[R = \{(x,y) : 0 \leq x \leq 1,\ x \leq y \leq 1\}
    = \{(x,y) : 0 \leq y \leq 1,\ 0 \leq x \leq y\}.\]
The first reads it by vertical slices, the second by horizontal ones. Taking the second and
applying \cref{thm:fubini},
\[I = \int_0^1\!\!\int_0^{y} e^{y^2}\,dx\,dy.\]

\textbf{3.} Now $e^{y^2}$ is constant in $x$, so the inner integral contributes a factor of $y$,
and that factor is exactly what the substitution $s := y^2$ needs:
\[I = \int_0^1 y\,e^{y^2}\,dy = \left[\frac{e^{y^2}}{2}\right]_0^1 = \frac{e-1}{2}.\]

\begin{remark}[What the swap actually bought]
\label{rem:order_swap_supplies_the_factor}
The reversal did not find a cleverer antiderivative. It produced a factor of $y$ that was not
there before, and $y\,e^{y^2}$ integrates where $e^{y^2}$ does not. That is the pattern worth
recognising: when the length of a slice depends on the variable that is causing the trouble,
swapping the order hands you that length as a factor in the integrand.
\end{remark}
\end{exercisesolution}

\begin{exercisesolution}[Solution to \cref{ex:ai_frullani}]
% Generator: Claude Opus 5 (High)
\textbf{1.} Both exponentials equal $1$ at $x = 0$, so the quotient is a difference quotient in
disguise. Expanding, $e^{-ax} - e^{-bx} = (b-a)x + O(x^2)$, and therefore
\[\lim_{x\to0^+}\frac{e^{-ax} - e^{-bx}}{x} = b - a.\]
Assigning that value at $x = 0$ makes the integrand continuous on $[0,\infty)$, and it decays
exponentially, so the improper integral converges.

\textbf{2.} Only the second exponential depends on $t$, and the $x$ that the chain rule produces
cancels the one in the denominator:
\[I'(t) = \int_0^\infty \frac{\partial}{\partial t}\left(\frac{e^{-ax} - e^{-tx}}{x}\right)dx
        = \int_0^\infty \frac{x\,e^{-tx}}{x}\,dx
        = \int_0^\infty e^{-tx}\,dx
        = \frac{1}{t}.\]
That cancellation is the whole trick. Differentiating in the parameter removes the $1/x$ which
made the integrand hard, and leaves an integral of one line.

\textbf{3.} At $t = a$ the two exponentials coincide and the integrand vanishes identically, so
$I(a) = 0$. Consequently, by the fundamental theorem of calculus,
\[\int_0^\infty \frac{e^{-ax} - e^{-bx}}{x}\,dx = I(b) = I(b) - I(a) = \int_a^b \frac{dt}{t}
  = \log\!\left(\frac{b}{a}\right).\]
That the answer sees only the ratio $b/a$ is now visible in the derivation rather than a
surprise, since $I'(t) = 1/t$ is scale-invariant.

\begin{remark}[The hypotheses of \cref{thm:feynmans_trick} do not literally cover this]
\label{rem:feynman_on_a_half_line}
As stated, \cref{thm:feynmans_trick} differentiates under an integral over a \emph{compact} set
$K$, and here the domain is $[0,\infty)$. The step is still legitimate, by the same reasoning
\cref{ex:feynman_dirichlet_integral} appeals to. Fix $0 < t_0 < t_1$. For $t \in [t_0,t_1]$ the
$t$-derivative of the integrand is $e^{-tx} \leq e^{-t_0x}$, a bound that is independent of $t$
and integrable on $[0,\infty)$, so the differentiation may be carried out on $[0,R]$ and the
limit $R \to \infty$ taken afterwards. The compact statement is the one the course proves; the
dominated version is what the applications actually use.
\end{remark}
\end{exercisesolution}

\begin{exercisesolution}[Proof of \cref{ex:volume_ellipsoid_change_variables}]
% Generator: Gemini 3.6 Flash (High)
% Correction: Claude Opus 5 (High) --- \operatorname{vol} -> \vol, J\Phi -> \Jac\Phi, ^T ->
% \transp, and \iiint -> \int as elsewhere in this chapter.
\begin{enumerate}[label=\textbf{(\alph*)}]
  \item Define $\Phi(u,v,w) := (au, bv, cw)\transp$, and write $(x,y,z) = \Phi(u,v,w)$. Then
    $(x/a)^2 + (y/b)^2 + (z/c)^2 = u^2+v^2+w^2$, so $(u,v,w)$ lies in the closed unit ball if and
    only if $(x,y,z) \in U$. Since $a,b,c > 0$, the map $\Phi$ is a linear bijection, and therefore
    $\Phi(\overline{B_1(0)}) = U$.

  \item The Jacobian matrix of $\Phi$ is the diagonal matrix
\[\Jac\Phi = \begin{pmatrix} a & 0 & 0 \\ 0 & b & 0 \\ 0 & 0 & c \end{pmatrix},\]
    constant in $(u,v,w)$, with determinant $\det \Jac\Phi = abc > 0$.

  \item By the change of variables theorem (\cref{thm:change_of_variables}),
\[\vol_3(U) = \int_U 1 \, d\vol_3 = \int_{\overline{B_1(0)}} \lvert\det \Jac\Phi\rvert \, d\vol_3 = abc \int_{\overline{B_1(0)}} 1 \, d\vol_3 .\]
The unit ball in $\mathbb{R}^3$ has volume $\tfrac{4}{3}\pi$, and therefore
\[\vol_3(U) = \frac{4}{3} \pi abc.\]
\end{enumerate}
\end{exercisesolution}

\begin{exercisesolution}[Proof of \cref{ex:cubic_shear_area}]
% Generator: Claude Opus 5 (High)
\begin{enumerate}[label=\textbf{(\alph*)}]
  \item Differentiating each component,
\[D_x\varphi = \Jac\varphi(x) = \begin{pmatrix} 3x_1^2 & -1 \\ 1 & 1 \end{pmatrix},
  \qquad \det \Jac\varphi(x) = 3x_1^2 + 1 \ge 1 > 0 .\]

  \item The determinant never vanishes, so by the inverse function theorem
    (\cref{thm:inverse_function_theorem}) $\varphi$ is a local diffeomorphism at every point; in
    particular $\varphi(\mathbb{R}^2)$ is open. It remains to prove global injectivity, and then
    $\varphi$ is a bijection onto its image which is locally a diffeomorphism, hence a
    diffeomorphism onto its image.

    Suppose $\varphi(x) = \varphi(y)$. The second components give $x_2 + x_1 = y_2 + y_1$, that is,
    $x_2 - y_2 = -(x_1 - y_1)$. Substituting this into the equality of first components,
    $x_1^3 - x_2 = y_1^3 - y_2$, yields
\[x_1^3 - y_1^3 = x_2 - y_2 = -(x_1 - y_1),
  \qquad\text{that is,}\qquad (x_1 - y_1)\big(x_1^2 + x_1y_1 + y_1^2 + 1\big) = 0 .\]
    The second factor is at least $1$, because $x_1^2 + x_1y_1 + y_1^2 = \big(x_1 +
    \tfrac{y_1}{2}\big)^2 + \tfrac34 y_1^2 \ge 0$. Hence $x_1 = y_1$, and then $x_2 = y_2$.

  \item By part \textbf{(a)} the Jacobian determinant is positive, so the change of variables
    theorem (\cref{thm:change_of_variables}) gives
\[\vol_2(\varphi(Q)) = \int_Q \lvert\det \Jac\varphi\rvert \, dx
  = \int_0^1\!\!\int_0^1 \big(3x_1^2 + 1\big) \, dx_2\,dx_1
  = \int_0^1 \big(3x_1^2+1\big)\,dx_1 = 1 + 1 = 2 .\]
\end{enumerate}
\begin{remark}
The map is not linear, so the image of the square is not a parallelogram and there is no single
factor by which areas scale. The change of variables formula handles that by scaling
\emph{pointwise} and integrating the result, which is the whole reason the determinant appears
inside the integral rather than outside it.
\end{remark}
\end{exercisesolution}

\begin{exercisesolution}[Solution of \cref{ex:gaussian_second_moment_polar}]
% Generator: Claude Opus 5 (High)
Following the hint, the substitution $(x,y) \mapsto (y,x)$ is a linear change of variables with
Jacobian determinant $-1$, and it carries the integrand $x^2e^{-\alpha(x^2+y^2)}$ to
$y^2e^{-\alpha(x^2+y^2)}$. Hence the two integrals agree, and adding them gives
\[2I_\alpha = \int_{\mathbb{R}^2} (x^2+y^2)\,e^{-\alpha(x^2+y^2)} \, dx\,dy .\]
The integrand is now radial, which is what makes polar coordinates pay. Writing $x = r\cos\theta$,
$y = r\sin\theta$, with Jacobian determinant $r$,
\[2I_\alpha = \int_0^{2\pi}\!\!\int_0^\infty r^2 e^{-\alpha r^2}\, r\,dr\,d\theta
  = 2\pi \int_0^\infty r^3 e^{-\alpha r^2}\,dr .\]
Substituting $u := r^2$, so that $du = 2r\,dr$, turns the remaining integral into one that can be
done by parts:
\[\int_0^\infty r^3 e^{-\alpha r^2}\,dr = \frac12 \int_0^\infty u e^{-\alpha u}\,du
  = \frac{1}{2\alpha^2}.\]
Therefore $2I_\alpha = 2\pi \cdot \frac{1}{2\alpha^2} = \frac{\pi}{\alpha^2}$, and
\[I_\alpha = \frac{\pi}{2\alpha^2}.\]
\begin{remark}
Two sanity checks come free. Scaling $(x,y) \mapsto (x,y)/\sqrt\alpha$ predicts $I_\alpha =
\alpha^{-2} I_1$, which the answer satisfies; and $I_\alpha$ should equal
$\tfrac12 \cdot \tfrac{1}{2\alpha}\cdot\tfrac{\pi}{\alpha}$ read off from the one-dimensional
Gaussian moments, which it does. The symmetry trick in the first line is the same one that
evaluates $\int_{\mathbb{R}} e^{-x^2}dx$ by squaring it: an integral that resists in Cartesian
coordinates becomes radial once it is paired with its own mirror image.
\end{remark}
\end{exercisesolution}



